% !TeX root = main.tex
\documentclass[11pt]{article}
% Shared preamble of the Guitar Virtual Amp study documents.
% Included by docs/papers/*/main.tex with % Shared preamble of the Guitar Virtual Amp study documents.
% Included by docs/papers/*/main.tex with % Shared preamble of the Guitar Virtual Amp study documents.
% Included by docs/papers/*/main.tex with \input{../common/preamble}.
%
% These are working documents, not submissions: one column, generous
% explanations, and explicit labels for where each claim comes from.

\usepackage[utf8]{inputenc}
\usepackage[T1]{fontenc}
\usepackage{lmodern}
\usepackage[a4paper,margin=2.6cm]{geometry}
\usepackage{microtype}
\usepackage{parskip}
\usepackage{graphicx}
\usepackage{tikz}
\usetikzlibrary{positioning,arrows.meta,fit,backgrounds,calc,shapes.geometric,decorations.pathreplacing}
\usepackage{pgfplots}
\pgfplotsset{compat=1.17}
\usepackage{booktabs}
\usepackage{tabularx}
\usepackage{longtable}
\usepackage{makecell}
\usepackage{multirow}
\usepackage{array}
\usepackage{amsmath,amssymb}
\usepackage[table]{xcolor}
\usepackage{enumitem}
\setlist{itemsep=2pt,topsep=4pt}
\usepackage{caption}
\captionsetup{font=small,labelfont=bf,margin=0.5cm}
\usepackage[numbers,sort&compress]{natbib}
\usepackage[skins,breakable]{tcolorbox}
\usepackage{fancyhdr}
\usepackage[colorlinks=true,linkcolor=blue!50!black,citecolor=green!40!black,urlcolor=blue!60!black]{hyperref}

\setcounter{tocdepth}{2}
\setcounter{secnumdepth}{3}

% ---------------------------------------------------------------- colours
\definecolor{keyc}{HTML}{1F5FAD}
\definecolor{appc}{HTML}{2E7D32}
\definecolor{warnc}{HTML}{C25E00}
\definecolor{examc}{HTML}{555555}
\definecolor{nextc}{HTML}{6A1B9A}

% ---------------------------------------------------------------- boxes
% Intuition before the details.
\newtcolorbox{keyidea}[1][]{enhanced,breakable,colback=keyc!5,colframe=keyc,
  fonttitle=\bfseries,title={Key idea#1},left=6pt,right=6pt,top=4pt,bottom=4pt}
% What Guitar Virtual Amp does, with code pointers.
\newtcolorbox{inapp}[1][]{enhanced,breakable,colback=appc!5,colframe=appc,
  fonttitle=\bfseries,title={In Guitar Virtual Amp#1},left=6pt,right=6pt,top=4pt,bottom=4pt}
% Pitfalls, caveats, things that are easy to misread.
\newtcolorbox{caution}[1][]{enhanced,breakable,colback=warnc!5,colframe=warnc,
  fonttitle=\bfseries,title={Caution#1},left=6pt,right=6pt,top=4pt,bottom=4pt}
% A small calculation done by hand.
\newtcolorbox{worked}[1][]{enhanced,breakable,colback=examc!4,colframe=examc,
  fonttitle=\bfseries,title={Worked example#1},left=6pt,right=6pt,top=4pt,bottom=4pt}
% Planned or in-progress work.
\newtcolorbox{nextsteps}[1][]{enhanced,breakable,colback=nextc!4,colframe=nextc,
  fonttitle=\bfseries,title={Next steps#1},left=6pt,right=6pt,top=4pt,bottom=4pt}

% ---------------------------------------------------------------- evidence labels
% Every number or claim about a method carries one of these.
\DeclareRobustCommand{\evtag}[2]{{\setlength{\fboxsep}{1.3pt}%
  \colorbox{#1!15}{\textcolor{#1!70!black}{\scriptsize\sffamily\bfseries #2}}}}
\DeclareRobustCommand{\measured}{\evtag{appc}{MEASURED}}  % run by us, in this repository
\DeclareRobustCommand{\reported}{\evtag{keyc}{REPORTED}}  % stated by the authors
\DeclareRobustCommand{\expected}{\evtag{warnc}{EXPECTED}} % hypothesis, not yet measured
% Status of a component of the app.
\DeclareRobustCommand{\done}{\evtag{appc}{BUILT}}
\DeclareRobustCommand{\wip}{\evtag{warnc}{IN PROGRESS}}
\DeclareRobustCommand{\planned}{\evtag{nextc}{PLANNED}}

% ---------------------------------------------------------------- shorthands
% Underscores and slashes need no escaping inside these.
\newcommand{\code}[1]{{\small\nolinkurl{#1}}}
\newcommand{\file}[1]{{\small\nolinkurl{#1}}}
% Shell commands and flags with spaces (\code drops spaces).
\newcommand{\cmd}[1]{\texttt{\small\detokenize{#1}}}
\newcommand{\ms}{\,ms}
\newcommand{\dB}{\,dB}
\newcommand{\Hz}{\,Hz}
\newcommand{\kHz}{\,kHz}
\newcommand{\term}[1]{\emph{#1}}

\pagestyle{fancy}
\fancyhf{}
\renewcommand{\headrulewidth}{0.3pt}
\fancyfoot[C]{\thepage}
.
%
% These are working documents, not submissions: one column, generous
% explanations, and explicit labels for where each claim comes from.

\usepackage[utf8]{inputenc}
\usepackage[T1]{fontenc}
\usepackage{lmodern}
\usepackage[a4paper,margin=2.6cm]{geometry}
\usepackage{microtype}
\usepackage{parskip}
\usepackage{graphicx}
\usepackage{tikz}
\usetikzlibrary{positioning,arrows.meta,fit,backgrounds,calc,shapes.geometric,decorations.pathreplacing}
\usepackage{pgfplots}
\pgfplotsset{compat=1.17}
\usepackage{booktabs}
\usepackage{tabularx}
\usepackage{longtable}
\usepackage{makecell}
\usepackage{multirow}
\usepackage{array}
\usepackage{amsmath,amssymb}
\usepackage[table]{xcolor}
\usepackage{enumitem}
\setlist{itemsep=2pt,topsep=4pt}
\usepackage{caption}
\captionsetup{font=small,labelfont=bf,margin=0.5cm}
\usepackage[numbers,sort&compress]{natbib}
\usepackage[skins,breakable]{tcolorbox}
\usepackage{fancyhdr}
\usepackage[colorlinks=true,linkcolor=blue!50!black,citecolor=green!40!black,urlcolor=blue!60!black]{hyperref}

\setcounter{tocdepth}{2}
\setcounter{secnumdepth}{3}

% ---------------------------------------------------------------- colours
\definecolor{keyc}{HTML}{1F5FAD}
\definecolor{appc}{HTML}{2E7D32}
\definecolor{warnc}{HTML}{C25E00}
\definecolor{examc}{HTML}{555555}
\definecolor{nextc}{HTML}{6A1B9A}

% ---------------------------------------------------------------- boxes
% Intuition before the details.
\newtcolorbox{keyidea}[1][]{enhanced,breakable,colback=keyc!5,colframe=keyc,
  fonttitle=\bfseries,title={Key idea#1},left=6pt,right=6pt,top=4pt,bottom=4pt}
% What Guitar Virtual Amp does, with code pointers.
\newtcolorbox{inapp}[1][]{enhanced,breakable,colback=appc!5,colframe=appc,
  fonttitle=\bfseries,title={In Guitar Virtual Amp#1},left=6pt,right=6pt,top=4pt,bottom=4pt}
% Pitfalls, caveats, things that are easy to misread.
\newtcolorbox{caution}[1][]{enhanced,breakable,colback=warnc!5,colframe=warnc,
  fonttitle=\bfseries,title={Caution#1},left=6pt,right=6pt,top=4pt,bottom=4pt}
% A small calculation done by hand.
\newtcolorbox{worked}[1][]{enhanced,breakable,colback=examc!4,colframe=examc,
  fonttitle=\bfseries,title={Worked example#1},left=6pt,right=6pt,top=4pt,bottom=4pt}
% Planned or in-progress work.
\newtcolorbox{nextsteps}[1][]{enhanced,breakable,colback=nextc!4,colframe=nextc,
  fonttitle=\bfseries,title={Next steps#1},left=6pt,right=6pt,top=4pt,bottom=4pt}

% ---------------------------------------------------------------- evidence labels
% Every number or claim about a method carries one of these.
\DeclareRobustCommand{\evtag}[2]{{\setlength{\fboxsep}{1.3pt}%
  \colorbox{#1!15}{\textcolor{#1!70!black}{\scriptsize\sffamily\bfseries #2}}}}
\DeclareRobustCommand{\measured}{\evtag{appc}{MEASURED}}  % run by us, in this repository
\DeclareRobustCommand{\reported}{\evtag{keyc}{REPORTED}}  % stated by the authors
\DeclareRobustCommand{\expected}{\evtag{warnc}{EXPECTED}} % hypothesis, not yet measured
% Status of a component of the app.
\DeclareRobustCommand{\done}{\evtag{appc}{BUILT}}
\DeclareRobustCommand{\wip}{\evtag{warnc}{IN PROGRESS}}
\DeclareRobustCommand{\planned}{\evtag{nextc}{PLANNED}}

% ---------------------------------------------------------------- shorthands
% Underscores and slashes need no escaping inside these.
\newcommand{\code}[1]{{\small\nolinkurl{#1}}}
\newcommand{\file}[1]{{\small\nolinkurl{#1}}}
% Shell commands and flags with spaces (\code drops spaces).
\newcommand{\cmd}[1]{\texttt{\small\detokenize{#1}}}
\newcommand{\ms}{\,ms}
\newcommand{\dB}{\,dB}
\newcommand{\Hz}{\,Hz}
\newcommand{\kHz}{\,kHz}
\newcommand{\term}[1]{\emph{#1}}

\pagestyle{fancy}
\fancyhf{}
\renewcommand{\headrulewidth}{0.3pt}
\fancyfoot[C]{\thepage}
.
%
% These are working documents, not submissions: one column, generous
% explanations, and explicit labels for where each claim comes from.

\usepackage[utf8]{inputenc}
\usepackage[T1]{fontenc}
\usepackage{lmodern}
\usepackage[a4paper,margin=2.6cm]{geometry}
\usepackage{microtype}
\usepackage{parskip}
\usepackage{graphicx}
\usepackage{tikz}
\usetikzlibrary{positioning,arrows.meta,fit,backgrounds,calc,shapes.geometric,decorations.pathreplacing}
\usepackage{pgfplots}
\pgfplotsset{compat=1.17}
\usepackage{booktabs}
\usepackage{tabularx}
\usepackage{longtable}
\usepackage{makecell}
\usepackage{multirow}
\usepackage{array}
\usepackage{amsmath,amssymb}
\usepackage[table]{xcolor}
\usepackage{enumitem}
\setlist{itemsep=2pt,topsep=4pt}
\usepackage{caption}
\captionsetup{font=small,labelfont=bf,margin=0.5cm}
\usepackage[numbers,sort&compress]{natbib}
\usepackage[skins,breakable]{tcolorbox}
\usepackage{fancyhdr}
\usepackage[colorlinks=true,linkcolor=blue!50!black,citecolor=green!40!black,urlcolor=blue!60!black]{hyperref}

\setcounter{tocdepth}{2}
\setcounter{secnumdepth}{3}

% ---------------------------------------------------------------- colours
\definecolor{keyc}{HTML}{1F5FAD}
\definecolor{appc}{HTML}{2E7D32}
\definecolor{warnc}{HTML}{C25E00}
\definecolor{examc}{HTML}{555555}
\definecolor{nextc}{HTML}{6A1B9A}

% ---------------------------------------------------------------- boxes
% Intuition before the details.
\newtcolorbox{keyidea}[1][]{enhanced,breakable,colback=keyc!5,colframe=keyc,
  fonttitle=\bfseries,title={Key idea#1},left=6pt,right=6pt,top=4pt,bottom=4pt}
% What Guitar Virtual Amp does, with code pointers.
\newtcolorbox{inapp}[1][]{enhanced,breakable,colback=appc!5,colframe=appc,
  fonttitle=\bfseries,title={In Guitar Virtual Amp#1},left=6pt,right=6pt,top=4pt,bottom=4pt}
% Pitfalls, caveats, things that are easy to misread.
\newtcolorbox{caution}[1][]{enhanced,breakable,colback=warnc!5,colframe=warnc,
  fonttitle=\bfseries,title={Caution#1},left=6pt,right=6pt,top=4pt,bottom=4pt}
% A small calculation done by hand.
\newtcolorbox{worked}[1][]{enhanced,breakable,colback=examc!4,colframe=examc,
  fonttitle=\bfseries,title={Worked example#1},left=6pt,right=6pt,top=4pt,bottom=4pt}
% Planned or in-progress work.
\newtcolorbox{nextsteps}[1][]{enhanced,breakable,colback=nextc!4,colframe=nextc,
  fonttitle=\bfseries,title={Next steps#1},left=6pt,right=6pt,top=4pt,bottom=4pt}

% ---------------------------------------------------------------- evidence labels
% Every number or claim about a method carries one of these.
\DeclareRobustCommand{\evtag}[2]{{\setlength{\fboxsep}{1.3pt}%
  \colorbox{#1!15}{\textcolor{#1!70!black}{\scriptsize\sffamily\bfseries #2}}}}
\DeclareRobustCommand{\measured}{\evtag{appc}{MEASURED}}  % run by us, in this repository
\DeclareRobustCommand{\reported}{\evtag{keyc}{REPORTED}}  % stated by the authors
\DeclareRobustCommand{\expected}{\evtag{warnc}{EXPECTED}} % hypothesis, not yet measured
% Status of a component of the app.
\DeclareRobustCommand{\done}{\evtag{appc}{BUILT}}
\DeclareRobustCommand{\wip}{\evtag{warnc}{IN PROGRESS}}
\DeclareRobustCommand{\planned}{\evtag{nextc}{PLANNED}}

% ---------------------------------------------------------------- shorthands
% Underscores and slashes need no escaping inside these.
\newcommand{\code}[1]{{\small\nolinkurl{#1}}}
\newcommand{\file}[1]{{\small\nolinkurl{#1}}}
% Shell commands and flags with spaces (\code drops spaces).
\newcommand{\cmd}[1]{\texttt{\small\detokenize{#1}}}
\newcommand{\ms}{\,ms}
\newcommand{\dB}{\,dB}
\newcommand{\Hz}{\,Hz}
\newcommand{\kHz}{\,kHz}
\newcommand{\term}[1]{\emph{#1}}

\pagestyle{fancy}
\fancyhf{}
\renewcommand{\headrulewidth}{0.3pt}
\fancyfoot[C]{\thepage}


\fancyhead[L]{\small Tone estimation notes}
\fancyhead[R]{\small\nouppercase{\leftmark}}

\title{Estimating the Guitar Tone of a Recording\\[4pt]
\large State of the art, how it works, and what Guitar Virtual Amp does}
\author{Breno de Angelo\\\small Guitar Virtual Amp project notes}
\date{Snapshot of 25 November 2026}

\begin{document}
\maketitle

\begin{abstract}
\noindent
A guitarist learning a song wants its notes and its sound: the amplifier and how
hard it is driven, the cabinet, the effects. This document explains how a
computer can estimate that sound (the \emph{tone}) from a recording and turn it
into settings of Guitar Virtual Amp's own pedals, what the research literature
has achieved, and what the app does today. As in the companion document on
tablature, every term is defined before use, numbers are labelled
\measured{}, \reported{} or \expected{}, and parts of the app \done{}, \wip{} or
\planned. The short version: recognising which effects are present works on
synthetic, isolated guitar, and searching a chain's settings by rendering it
(\emph{analysis by synthesis}) beats predicting them; but no published method
estimates an editable amplifier, cabinet and effects chain from a guitar
separated out of a commercial song. The app now has a first version of exactly
that: a gradient-free search (CMA-ES) over its own compressor, distortion, amp,
delay and reverb, driven by a stand-in dry signal synthesised from the
transcribed notes, followed by a matching filter. On 12 EGDB test items it beats
the best factory preset on every target group, but most of the remaining error
comes from the stand-in dry signal, not from the search.
\end{abstract}

\tableofcontents
\clearpage

% !TeX root = ../main.tex
\section*{How to read this document}
\addcontentsline{toc}{section}{How to read this document}

\paragraph{Purpose.} Coding agents are building a tone estimator for Guitar
Virtual Amp (\file{guitar_core/src/tone/}). This document explains the problem,
the research it builds on, the design that was chosen, exactly what is
implemented, and how well it works, so the owner can follow and steer that
work. It replaces an earlier conference-style literature review and keeps all of
its references.

\paragraph{Structure.}
\begin{itemize}
\item \textbf{Part I, Background} (Sections~\ref{sec:problem}--\ref{sec:search}):
  what a tone is, the signal processing behind a guitar rig, how two sounds are
  compared, and how a search finds settings. Read this first if any of
  ``impulse response'', ``minimum phase'' or ``CMA-ES'' is unfamiliar.
\item \textbf{Part II, The literature} (Sections~\ref{sec:datasets}--\ref{sec:comparison}):
  datasets, effect recognition, amplifier modelling, working without the dry
  signal, style transfer and sound matching, commercial products, and a
  comparison with failure modes.
\item \textbf{Part III, Guitar Virtual Amp} (Sections~\ref{sec:design}--\ref{sec:next}):
  the design, the implementation component by component (with the formulas the
  code uses), the benchmark and its results, and the next steps.
\item \textbf{Appendices}: glossary; code map and commands.
\end{itemize}

\paragraph{Labels.} Coloured boxes: \emph{Key idea} (intuition), \emph{In Guitar
Virtual Amp} (what the code does), \emph{Caution} (pitfalls), \emph{Worked
example}, \emph{Next steps}. Inline tags: \measured{} run by us in this
repository; \reported{} stated by a paper's authors; \expected{} a hypothesis.
For the app: \done{} in the committed code, \wip{} in the working tree or being
changed, \planned{} only in the roadmap.

\paragraph{Snapshot and moving parts.} Written on 25 November 2026, updated on
26 November. The tone module was committed in \code{7db5015} (search, pluck
stand-in, match IR), \code{25f53ea} (measured delay, separated-stem targets, the
tone stage of the transcription pipeline) and \code{9939922} (preset prior, two
saved presets). The 26 November update adds a breakdown of where the error comes
from and a colour calibration of the stand-in, measured in the benchmark but not
yet used by the app (Section~\ref{sec:results-colour}). The results in
Section~\ref{sec:results} are from the committed defaults.

\paragraph{Related files.} \file{ROADMAP.md} (``Tone from the recording''),
\file{CLAUDE.md} (section ``Tone from a recording''), and the companion
document on tablature, \file{docs/papers/tablature-transcription/}, whose
separation and note transcription this work reuses~\cite{deangelo2026tabreview}.

\subsection*{The whole story in one table}
\addcontentsline{toc}{subsection}{The whole story in one table}

\begin{table}[h]
\centering
\caption{Summary. Distances are in dB of log-band difference (lower is better,
Section~\ref{sec:app-distance}).}
\label{tab:summary}
\small
\begin{tabularx}{\textwidth}{@{}>{\raggedright}p{2.6cm}>{\raggedright}p{4.4cm}X@{}}
\toprule
Question & Literature & Guitar Virtual Amp \\
\midrule
Which effects are present? & Close to solved on synthetic isolated guitar (set F1 0.987 on 3-effect chains) \reported & Not estimated separately: the search switches pedals off \\
Knob values of a known chain? & Regression MAE 0.02--0.05 on a unit scale for known plugins \reported; search by re-rendering beats prediction (wet SI-SDR 18.2 $\to$ 23.1\dB) \reported & CMA-ES search through the app's own renderer \done \\
Without the dry signal? & Unpaired amp models score 77--84 MUSHRA \reported; effect removal; no notes-based resynthesis found & Dry signal synthesised from the transcribed notes (plucked-string model) \done; its colour calibrated on the player's DI measured \wip \\
On a stem from a commercial mix? & Only one measurement: effect retrieval halves on separated stems \reported & Separated-stem targets measured: search and preset + match tie (5.9 and 5.8\dB) \measured \\
Final colour (cabinet, mix EQ) & Spectral matching to an IR, after the gain is right (Tone Match) & Minimum-phase match IR in a Cabinet pedal \done \\
How good is it? & No benchmark for the full task & EGDB, 12 items: 5.6--6.5\dB{} against 5.8--6.9\dB{} for preset + match; 1.7--3.1\dB{} with the true DI; 5.0--5.9\dB{} with a calibrated stand-in \measured \\
\bottomrule
\end{tabularx}
\end{table}


\part{Background}
% !TeX root = ../main.tex
\section{The problem}
\label{sec:problem}

\subsection{What a ``tone'' is}

A recorded guitar part is the product of the notes and of the equipment that
shaped them (Figure~\ref{fig:rig}): the guitar and its pickups; often a
compressor or a drive pedal; an amplifier, whose \term{preamp} distorts, whose
\term{tone stack} (bass, middle, treble knobs) filters and whose \term{power
amp} adds more distortion; a speaker \term{cabinet} recorded by a microphone;
modulation (chorus, phaser, flanger), delay and reverb; and the mixing and
mastering of the whole song. Guitarists call the combined result the
\term{tone}.

\begin{figure}[h]
\centering
\resizebox{\textwidth}{!}{%
\begin{tikzpicture}[
  node distance=3mm,
  box/.style={draw, rounded corners=2pt, font=\small, align=center,
              minimum height=9mm, inner sep=3pt, fill=white},
  arr/.style={-{Stealth[length=2mm]}, semithick}]
  \node[box] (g) {Guitar,\\pickups};
  \node[box, right=of g] (p) {Pedals\\(comp., drive)};
  \node[box, right=of p] (a) {Amp: preamp,\\tone stack, power};
  \node[box, right=of a] (c) {Cabinet\\+ mic};
  \node[box, right=of c] (t) {Modulation,\\delay, reverb};
  \node[box, right=of t] (m) {Mix and\\master};
  \foreach \x/\y in {g/p,p/a,a/c,c/t,t/m} \draw[arr] (\x) -- (\y);
  \draw[decorate,decoration={brace,amplitude=4pt,mirror}] ([yshift=-2mm]g.south west) -- node[below=4pt,font=\scriptsize]{the app cannot change these, only compensate} ([yshift=-2mm]g.south east);
  \draw[decorate,decoration={brace,amplitude=4pt}] ([yshift=2mm]p.north west) -- node[above=4pt,font=\scriptsize]{what the app's pedals can reproduce} ([yshift=2mm]t.north east);
  \draw[decorate,decoration={brace,amplitude=4pt,mirror}] ([yshift=-2mm]m.south west) -- node[below=4pt,font=\scriptsize,align=center]{absorbed by the\\final EQ, partly} ([yshift=-2mm]m.south east);
\end{tikzpicture}}
\caption{A guitar signal chain. The \emph{dry} signal (DI) is the output of
the first box; a recording is the output of the last.}
\label{fig:rig}
\end{figure}

A player who learns a song wants this sound, and today finds it by ear, by
reading interviews about the artist's gear, or by buying a preset made for the
song. The goal here is to propose, automatically, a preset of Guitar Virtual
Amp's pedals that sounds like the guitar in a recording the user owns.

\subsection{An inverse problem}

\begin{keyidea}[: forward and inverse]
The \textbf{forward} problem is easy: given a dry signal $x$ and settings
$\theta$ of a chain $f$, compute the sound $y = f(x;\theta)$. The app does this
in real time. The \textbf{inverse} problem is hard: given only a recording
$y^\star$, find settings $\hat\theta$ such that $f(x;\hat\theta)$ sounds like
$y^\star$. We do not even know $x$, the dry signal that went into the original
chain.
\end{keyidea}

Formally, the estimate is
\[
  \hat\theta = \arg\min_{\theta}\ D\bigl(f(\tilde x;\theta),\ y^\star\bigr),
\]
where $D$ measures how different two sounds are (Section~\ref{sec:metrics}),
$\tilde x$ is the best dry signal we can get, and the minimisation is done by a
search (Section~\ref{sec:search}). Every part of this formula is a design
choice, and the literature differs on each.

\subsection{What makes it hard}

\begin{itemize}
\item \textbf{Nonlinear, time-varying, long memory.} Distortion is nonlinear,
  compression and modulation change over time, delay and reverb remember the
  past for seconds (Section~\ref{sec:dsp}).
\item \textbf{Many settings sound alike.} Amp gain, a drive pedal and a
  compressor can all add sustain and grit; different orders of the same
  effects can sound similar. The inverse problem has no unique answer, so
  ``recover the original settings'' is the wrong goal; ``sound alike'' is the
  right one.
\item \textbf{No dry signal.} Amplifier capture (how Kemper or Neural Amp
  Modeler work) needs the dry input and the processed output of the same
  performance, sample by sample. A song gives only the output.
\item \textbf{The mix.} The guitar must first be separated from drums, bass,
  keys and voice; separation leaves artefacts that look like effects, and
  bus compression and mastering act after the guitarist's chain.
\item \textbf{Different gear.} The target is not the original equipment but its
  sound on the equipment the player has: here, a mono chain of the app's 18
  pedals.
\end{itemize}

\subsection{Three settings, by what is available}

Methods differ most in what they require:
\begin{description}[style=nextline,leftmargin=1.2cm]
\item[Paired] The dry signal (direct input, \term{DI}) and the processed signal
  of the same performance, aligned sample by sample. This is how amplifiers are
  captured.
\item[Unpaired] A processed recording plus dry guitar from \emph{another}
  performance, maybe on another guitar.
\item[Wet only] Only the processed recording.
\end{description}
A song is wet only, and separated from a mix. The player's own guitar can
supply unpaired dry audio. The app adds a fourth option that the literature does
not use: a dry signal \emph{synthesised from the transcribed notes}, which is
aligned with the target but has the wrong guitar timbre
(Section~\ref{sec:impl-dry}).

\subsection{What the output should be}

A tone estimate can be (i) an \term{interpretable chain}: effect types with
parameter values that a player can load and adjust; (ii) a \term{neural model}
of the processing, such as an amplifier capture, which reproduces a sound but
exposes few controls; or (iii) only \term{processed audio}, from tone-transfer
models that re-render a recording in another tone. Only (i) and (ii) can be
used in a real-time effects app, and only (i) can be edited. The app aims for
(i): a preset of its own pedals.

% !TeX root = ../main.tex
\section{Signal processing primer}
\label{sec:dsp}

The pieces of a guitar rig fall into four kinds, and each kind calls for a
different estimation strategy. This section explains them.

\subsection{Linear filters: EQ, tone stack, cabinet}

A processor is \term{linear and time-invariant} (LTI) when (a) doubling the
input doubles the output, (b) the response to a sum of inputs is the sum of the
responses, and (c) it behaves the same at all times. Equalisers, the tone stack
of an amplifier (taken alone), a speaker cabinet with its microphone, and to a
good approximation a room, are LTI.

\begin{keyidea}[: an LTI system is fully described by its impulse response]
Feed an LTI system a single click (an \term{impulse}); the output is its
\term{impulse response} (IR) $h$. The output for any input $x$ is then the
\term{convolution}
\[
  y[n] = \sum_{k} h[k]\,x[n-k],
\]
each input sample triggering a scaled copy of the IR. Equivalently, in
frequency, each frequency is multiplied by the \term{frequency response}
$H(f)$, the Fourier transform of $h$. This is why a cabinet can be stored as a
\file{.wav} IR and why the app's Cabinet pedal is a convolution
(\file{backend/dsp/convolution.rs}).
\end{keyidea}

Two consequences matter for tone estimation.
\begin{itemize}
\item \textbf{The linear part can be measured from long-term spectra.} If
  $y = h * x$, the average spectrum of $y$ is the average spectrum of $x$
  times $|H(f)|^2$, so in decibels
  $L_y(f) = L_x(f) + 20\log_{10}|H(f)|$. The difference between two
  \term{long-term average spectra} is the EQ that turns one into the other.
  This is what Tone Match products and the app's final step do.
\item \textbf{Phase.} Many IRs have the same magnitude response $|H(f)|$. The
  \term{minimum-phase} one concentrates its energy as early as possible, so it
  adds the least delay and smearing~\cite{oppenheim2010dtsp}. When only a
  magnitude is known (as from spectra), the minimum-phase IR is the natural
  choice. It can be computed through the \term{cepstrum} (the inverse Fourier
  transform of the log magnitude): keeping only its causal half and
  transforming back gives the minimum-phase filter (Section~\ref{sec:impl-match}).
\end{itemize}

\subsection{Nonlinear processors: drive, distortion, the amplifier}

A distortion circuit \term{clips}: above some level, the output stops growing.
Clipping a sine wave adds harmonics (multiples of its frequency); clipping two
notes together adds sum and difference frequencies (\term{intermodulation}).
Two properties of nonlinear processors drive the design of the app's search.

\begin{keyidea}[: input level is a setting]
How much a clipper distorts depends on how loud its input is. The same amp at
the same gain sounds clean with a soft guitar and crunchy with a hot pickup. A
tone estimate must therefore choose the input level together with the gain; in
the app it is the first searched parameter (\code{input level (RMS)}).
\end{keyidea}

\begin{worked}[: order matters]
Put a treble boost \emph{before} a distortion: the high strings drive the
clipper harder, so they distort more and gain harmonics, and the result is
brighter and fizzier. Put the same boost \emph{after}: the distortion is the
same as without it and only the final colour is brighter. For an LTI chain, the
order of filters does not matter; around a nonlinearity it does. Corollary: a
linear filter \emph{after} all the clipping (the cabinet, mix EQ, and in the
app's amp model the tone stack) can be estimated separately at the end, which is
what the app does.
\end{worked}

\subsection{Dynamics: compressors}

A \term{compressor} lowers the gain when the signal exceeds a
\term{threshold}, by a \term{ratio} (4:1 means 4\dB{} above the threshold
becomes 1\dB), reacting with an \term{attack} time and recovering with a
\term{release} time. It changes the envelope of notes (less attack, more
sustain) more than their spectrum. It is nonlinear and time-varying, and it is
easy to confuse with the natural compression of a distorted amp.

\subsection{Time-varying effects: modulation}

Chorus, flanger, phaser and tremolo are driven by a \term{low-frequency
oscillator} (LFO, typically 0.1--10\Hz) that slowly moves a delay time, a
filter notch or the volume. Their sound depends on the LFO rate and depth. The
search landscape over an LFO rate is very rugged (many local minima): a small
change in rate shifts the LFO's phase completely after a few seconds.

\subsection{Long-memory effects: delay and reverb}

A \term{delay} adds copies of the signal after a fixed time (the \term{delay
time}), each copy quieter by the \term{feedback} factor. A \term{reverb} adds a
dense, decaying tail; its decay time (RT60, the time to fall by 60\dB) and
brightness are its main traits. Like LFO rates, delay times make rugged
landscapes: a delay of 300\ms{} and one of 310\ms{} produce echoes that land in
different places, so the comparison score does not change smoothly between
them. That is why the app \emph{measures} the delay time from the recording
instead of searching it (Section~\ref{sec:impl-delay}).

\subsection{Loudness}

Two renders that differ only in volume should count as the same tone, because
the mix level of a guitar says nothing about its rig. All comparisons in the app
first scale both signals to the same RMS (root-mean-square) level.

% !TeX root = ../main.tex
\section{Comparing two sounds}
\label{sec:metrics}

The distance $D$ in the estimation formula decides what ``sounds alike'' means.
It is used twice: as the \emph{objective} the search minimises, and as the
\emph{score} a benchmark reports. They need not be the same.

\subsection{Parameter metrics}

If the target was made with the same chain the estimator controls, its true
settings are known, and one can report effect-detection accuracy or F1, the
accuracy of the predicted order, or the mean absolute error (MAE) of knob values
on a 0--1 scale. These need ground truth and say little about the sound: because
different settings can sound alike, a large parameter error can come with a good
sound. Peladeau and Peeters found that training on an audio loss gives better
reconstructions than training on parameters, even when the parameter error is
larger~\cite{peladeau2024blind}. Parameter error should therefore be a
diagnostic, not the main criterion.

\subsection{Aligned audio metrics}

These compare a render $\hat y$ with the target $y$ sample by sample or frame by
frame, so they need the \emph{same performance} on both sides.

\begin{description}[style=nextline,leftmargin=1.2cm]
\item[Error-to-signal ratio (ESR)]
  $\displaystyle \mathrm{ESR} = \frac{\sum_n (y[n]-\hat y[n])^2}{\sum_n y[n]^2}$,
  usually after a pre-emphasis filter that weights high frequencies the way
  hearing does~\cite{wright2020perceptual}. The standard loss for training amp
  models; 0 is perfect. Very sensitive to phase.
\item[Multi-resolution STFT distance (MR-STFT)] The difference of magnitude
  spectrograms at several window sizes (e.g.\ 512, 1024, 2048), as a sum of a
  spectral-convergence term and the L1 distance of log
  magnitudes~\cite{yamamoto2020parallel,steinmetz2020auraloss}. Ignores phase;
  the several resolutions catch both fast transients and fine pitch detail.
\item[SI-SDR] Scale-invariant signal-to-distortion ratio: project $\hat y$ on
  $y$ to get the best-scaled target $\alpha y$, then
  $10\log_{10}\bigl(\|\alpha y\|^2/\|\hat y - \alpha y\|^2\bigr)$, in dB, higher
  is better~\cite{leroux2019sdr}.
\item[Log-band distance (the app's)] Mean absolute difference, in dB, between
  log energies in a set of frequency bands, over frames; defined exactly in
  Section~\ref{sec:app-distance}. Close in spirit to the log-magnitude part of
  MR-STFT.
\end{description}

\subsection{Alignment-free metrics}

These compare clips without requiring the same notes:
\begin{itemize}
\item \textbf{Effect embeddings}: networks trained to describe the processing
  independently of the content, such as AFx-Rep~\cite{steinmetz2024stito} and
  Fx-Encoder++~\cite{yeh2025fxencoder}, or general audio embeddings such as
  CLAP~\cite{wu2023laionclap}. Two clips are compared by the cosine distance of
  their embeddings.
\item \textbf{Distribution distances}: the Fréchet Audio Distance
  (FAD)~\cite{kilgour2019fad} fits a Gaussian to the embeddings of many clips of
  each set and measures how far apart the two Gaussians are. It is biased by
  sample size and depends on the embedding and the reference set; FAD$_\infty$
  extrapolates the bias away~\cite{gui2024adapting}, and the kernel variant KAD
  converges with fewer samples~\cite{chung2025kad}. A single song gives few clips,
  so per-song evaluation suits pairwise embedding distances better than FAD.
\end{itemize}

\subsection{Listening tests}

The reference is human judgement. \term{MUSHRA} (ITU-R BS.1534-3)~\cite{itu2015bs1534}
asks listeners to rate several versions against a reference on a 0--100 scale,
with a hidden copy of the reference and a degraded anchor as checks;
webMUSHRA~\cite{schoeffler2018webmushra} runs it in a browser. \term{Mean opinion
scores} (1--5) can be split into tone similarity and audio
quality~\cite{yeh2026egvae}. As a ceiling, listeners could barely tell Kemper
profiles from the original amplifiers ($d' = 0.34$, $N = 177$)~\cite{duvel2020confusingly}.

\subsection{Which metrics agree with listeners}

Two studies measured this for guitar devices (Table~\ref{tab:metrics}). In the
ToneTwist evaluation of 40 nonlinear devices (27 screened listeners), MR-STFT
had the strongest correlation with MUSHRA ratings among signal metrics and ESR
one of the weakest; among FAD embeddings, VGGish and AFx-Rep did best overall
and CLAP best on distortion; no metric was reliable for overdrive or compression
\reported~\cite{comunita2025differentiable}. Wright, Dal Santo and Schlecht
collected ratings for 590 emulation--reference pairs and found MFCC and deep
embedding distances predictive, MAE and ESR poor \reported~\cite{wright2026predicting}.
In unpaired amp modelling, ESR and listener ratings disagreed on a heavily
distorted tone~\cite{wright2023adversarial}. No metric has been validated for
delay, reverb or modulation on guitar.

\begin{table}[h]
\centering
\caption{Spearman correlation of objective metrics with MUSHRA ratings on 40
nonlinear guitar and bass devices~\cite{comunita2025differentiable}. Negative:
lower distance goes with higher rating, so more negative is better.}
\label{tab:metrics}
\small
\begin{tabular}{@{}lrlr@{}}
\toprule
Signal metric & $\rho$ & FAD embedding & $\rho$ \\
\midrule
MR-STFT & $-0.427$ & VGGish & $-0.430$ \\
MAPE & $-0.339$ & AFx-Rep & $-0.417$ \\
L1 & $-0.332$ & CLAP & $-0.406$ \\
MSE & $-0.314$ & PANN & $-0.303$ \\
ESR & $-0.311$ & & \\
\midrule
\multicolumn{4}{@{}l@{}}{High-distortion devices only: MR-STFT $-0.536$, FAD with CLAP $-0.613$.} \\
\bottomrule
\end{tabular}
\end{table}

\begin{caution}[: correlations are modest]
The best correlation in Table~\ref{tab:metrics} is about $-0.43$: objective
metrics explain only part of what listeners hear. A lower distance is evidence
of a better match, not proof. Listening tests are on the app's roadmap
(Section~\ref{sec:next}); the benchmark already writes WAV files for listening
(\code{GVA_TONE_OUT}).
\end{caution}

\subsection{Comparability of published results}

Most results use clean, isolated, plugin-rendered guitar, which is easier than
the commercial case. Many test sets share dry recordings or plugins with the
training set. And because settings can be non-unique, two estimates with very
different parameters may sound alike, which parameter metrics penalise and
listeners do not.

% !TeX root = ../main.tex
\section{Finding settings: prediction and search}
\label{sec:search}

\subsection{Two ways to invert a chain}

\begin{description}[style=nextline,leftmargin=1.2cm]
\item[Prediction] Train a network on many renders of the chain with random
  settings to output the settings from the audio. Fast at use time (one pass),
  but it only knows the kind of audio it was trained on, and it struggles when
  several settings sound alike (it averages them).
\item[Analysis by synthesis (search)] Propose settings, render the dry signal
  through the real chain, measure the distance to the target, and improve the
  proposal; repeat. Slow (hundreds of renders), but it uses the exact forward
  model and needs no training.
\end{description}
The literature agrees on combining them: a prediction gives starting points, a
search refines them (Sections~\ref{sec:recognition} and~\ref{sec:matching}).
The app currently implements the search, starting from neutral settings and
from the factory presets.

\subsection{Gradient or no gradient}

If the chain is written in a differentiable framework (DDSP, PyTorch), the
gradient of the distance with respect to every setting is available and
gradient descent can be used. The app's chain is Rust code working on 16-bit
samples, with switches (pedal on or off) and discrete choices (amp voicing): it
has no gradients. So the search must be \term{gradient-free}: it may only
evaluate the distance at chosen points. That is also the situation of the
commercial plugins ST-ITO and Okita and Katayose search over
(Section~\ref{sec:recognition}). On non-convex problems gradient-free search can
even win: in synthesizer matching, CMA-ES reached a loss of 2.09 where Adam
(a gradient method) reached 4.60 \reported~\cite{bogdan2026instrumental}.

\subsection{CMA-ES, step by step}

\begin{keyidea}[: a cloud that moves, stretches and shrinks]
CMA-ES (Covariance Matrix Adaptation Evolution Strategy)~\cite{hansen2001cmaes}
keeps a Gaussian ``cloud'' over the search space: a \emph{mean} (the current
best guess), a \emph{step size} $\sigma$ (how wide the cloud is) and a
\emph{covariance} (its shape: which directions are stretched). Each generation
it samples candidates from the cloud, evaluates them, moves the mean toward the
best ones, stretches the cloud along directions that worked, and widens or
narrows it depending on whether the steps have been consistently in one
direction (far from the optimum: widen) or random (near it: narrow).
\end{keyidea}

One generation, with $n$ parameters, population $\lambda$ and the best
$\mu = \lambda/2$ kept:
\begin{enumerate}
\item \textbf{Sample} $\lambda$ candidates $x_k = m + \sigma\, C^{1/2} z_k$,
  with $z_k$ standard normal.
\item \textbf{Evaluate} each candidate (here: render and compare). This is the
  expensive part and runs in parallel.
\item \textbf{Rank} them and form the new mean as a weighted average of the
  best $\mu$, with larger weights for better ranks:
  $m' = \sum_{i=1}^{\mu} w_i\, x_{(i)}$.
\item \textbf{Update the shape} $C$ from the steps of the best candidates
  (rank-$\mu$ update) and from the accumulated path of the mean (rank-one
  update), so directions of progress grow.
\item \textbf{Update the step size}: compare the length of the accumulated path
  of the mean with the length expected of a random walk; longer means the
  search is moving consistently, so $\sigma$ grows; shorter means it is
  circling the optimum, so $\sigma$ shrinks.
\end{enumerate}
Only the \emph{ranks} of the candidates matter, not the distance values, which
makes CMA-ES robust to the scale and noise of the objective. The
\term{separable} variant~\cite{ros2008sepcmaes} keeps only the diagonal of $C$
(one width per parameter, no correlations): cheaper and faster to adapt, and
adequate for about ten parameters. The app uses it (\file{tone/cmaes.rs}).

\begin{worked}[: budget]
The app runs three searches (one per amp voicing), each with a population of
16 for 30 generations: $3 \times (1 + 16\cdot30) = 1443$ renders, plus one
render per factory preset for the baseline. Each render processes a 12\,s
excerpt through the chain. In the committed benchmark one estimate took about
15\,s on 28 threads \measured.
\end{worked}

\subsection{Pitfalls of search}

\begin{itemize}
\item \textbf{Rugged landscapes}: LFO rates and delay times
  (Section~\ref{sec:dsp}). Okita and Katayose had to fix the chorus rate because
  its landscape was too rugged \reported~\cite{okita2026audio}. Remedy: measure
  those parameters from the signal instead.
\item \textbf{Discrete choices} (amp voicing, which capture, which IR) are not
  continuous: search each one separately, or treat them as candidates.
\item \textbf{The objective can be gamed}: TAPE showed a search driving an
  embedding distance down while the spectral error rose above that of the
  unprocessed input \reported~\cite{long2026tape}. Use an objective close to
  what listeners hear, and check the result with a different score.
\item \textbf{Ambiguity}: many settings reach the same distance. A
  \term{prior} over realistic presets (a probability distribution fitted on
  real presets) turns the search into a maximum-a-posteriori estimate and
  reduced parameter error by up to a third in Yu et al.\ \reported~\cite{yu2025improving}.
\end{itemize}


\part{The literature}
% !TeX root = ../main.tex
\section{Datasets and licences}
\label{sec:datasets}

To train or test a tone estimator we need processed guitar audio, ideally with
the dry signal of the same performance (so the true chain can be re-applied) and
labels saying which equipment and settings produced it. Table~\ref{tab:datasets}
lists what exists. The sets fall into four groups.

\paragraph{Effect sets with dry and processed notes.} IDMT-SMT-Audio-Effects:
55\,044 two-second notes and chords of two guitars and two basses through ten
software effects and a no-effect class~\cite{stein2010aes}. GUITAR-FX-DIST
re-renders its dry notes through 13 overdrive, distortion and fuzz plugins with
continuous gain and tone (about 250\,h)~\cite{comunita2021guitar}. EGFxSet
records every note of a Stratocaster on five pickup positions through 12
hardware pedals at one setting each~\cite{pedroza2022egfxset}. Hinrichs et al.\
place processed guitar inside mixes with sample-based bass, keys and drums, the
only set with effects estimated in a mix~\cite{hinrichs2022convolutional} (its
release could not be confirmed). Rossi et al.\ record 400 DI improvisations on
four guitars and apply five effects in 32 chains with fixed or random parameters
and order~\cite{rossi2025automatic}.

\paragraph{Amplifier renders of DI performances.} EGDB: 118 minutes of DI from a
hexaphonic pickup, rendered through five amplifier simulations~\cite{chen2022egdb}.
EGDB-PG renders the same DI through 256 BIAS FX\,2 head and cabinet combinations
labelled low gain, crunch or high gain (514\,h)~\cite{chen2026toward}. GOAT
pairs 5.9\,h of DI with tablature and renders each item five times through
randomly drawn public NAM captures, a cabinet IR and an optional room
reverb~\cite{loth2025goat}. Guitar-TECHS records DI, a microphone on three real
amplifiers at flat EQ, and room microphones~\cite{pedroza2025guitartechs};
EGSet12 has 12 performances through a miked amplifier, without
DI~\cite{pedroza2024egset12}.

\paragraph{Device sets with knob labels.} ToneTwist AFx: dry and processed pairs
from over 40 compressors, overdrives, distortions, fuzzes and amplifiers, some
densely sampled over the knobs~\cite{comunita2025differentiable}. Amp-Space
sampled real devices with servo-driven knobs (about 50\,h real, 500\,h
simulated)~\cite{naradowsky2021ampspace}. Open-Amp is a generator that renders
dry audio through 160 crowd-sourced LSTM captures on the fly~\cite{wright2025openamp}.

\paragraph{Real multitracks.} MoisesDB splits guitar stems into acoustic, clean
electric and distorted electric (240 songs)~\cite{pereira2023moisesdb};
MedleyDB has the same classes and some raw tracks~\cite{bittner2014medleydb};
MUSDB18 keeps guitar inside ``other''~\cite{rafii2017musdb18}. None has
equipment or settings labels, and we found no open set that pairs commercial
guitar stems with verified equipment.

\begin{table}[h]
\centering
\caption{Datasets with processed guitar audio. ``DI'': the dry signal of the
same performance is included. Licences as stated on the hosting record on
2026-11-25; NC = non-commercial, ND = no derivatives.}
\label{tab:datasets}
\footnotesize
\setlength{\tabcolsep}{3pt}
\resizebox{\textwidth}{!}{%
\begin{tabular}{@{}lllllll@{}}
\toprule
Dataset & Audio & DI & Labels & Size & Processing & Licence \\
\midrule
IDMT-SMT-Audio-Effects~\cite{stein2010aes} & notes, chords & yes & effect class & 55k notes & 10 plugins & CC BY-NC-ND \\
GUITAR-FX-DIST~\cite{comunita2021guitar} & IDMT notes & yes & plugin, gain, tone & $\approx$250\,h & 13 drive plugins & CC BY (dry NC-ND) \\
EGFxSet~\cite{pedroza2022egfxset} & single notes & yes & pedal, 1 setting & 8,970 clips & 12 hardware pedals & CC BY \\
Rossi et al.~\cite{rossi2025automatic} & improvisations & yes & chain, order & 400 DI & 5 effects & CC BY \\
EGDB~\cite{chen2022egdb} & tablatures & yes & notes & 118\,min $\times$ 6 & 5 amp sims & not stated \\
EGDB-PG~\cite{chen2026toward} & EGDB DI & yes & preset, gain class & 514\,h & 256 head/cab presets & CC BY \\
GOAT~\cite{loth2025goat} & performances & yes & tablature & 5.9\,h DI & NAM captures + IR & CC BY-NC, on request \\
Guitar-TECHS~\cite{pedroza2025guitartechs} & techniques, excerpts & yes & MIDI notes & $>$5\,h & 3 real amps, miked & CC BY \\
EGSet12~\cite{pedroza2024egset12} & solos & no & notes, strings & 380\,s & 1 real amp, miked & CC BY \\
ToneTwist AFx~\cite{comunita2025differentiable} & guitar, bass & yes & device, knobs & 40+ devices & real hardware & mostly CC BY-NC \\
Amp-Space~\cite{naradowsky2021ampspace} & various & yes & device, knobs & 50\,h real & servo-driven devices & not stated \\
MoisesDB~\cite{pereira2023moisesdb} & songs & no & clean / distorted & 240 songs & real productions & CC BY-NC-SA \\
\bottomrule
\end{tabular}}
\end{table}

\subsection{Captures, impulse responses and their terms}

Libraries of amp captures and cabinet IRs are the natural source of amplifier
variety, but their terms matter. TONE3000 forbids automated or bulk downloading
and data mining in its site terms~\cite{tone30002025terms}; its API terms forbid
bulk download, caching or preloading of the catalogue and require that models are
downloaded only when a user asks for them~\cite{tone30002026apiterms}. Neither
grants a right to train on the catalogue, so GOAT-style renders through thousands
of public captures are a research resource we cannot reproduce from TONE3000.
Common free cabinet IRs (Kalthallen, OwnHammer) forbid redistribution and
modification~\cite{kalthallenfree,ownhammereula}. NAM's standard re-amping
signal (about three minutes at 48\kHz) is a fixed probe for making captures
in-house~\cite{atkinsonndnam}. The sets that allow redistribution and contain real
playing with a DI are EGDB-PG's renders, Guitar-TECHS, EGFxSet, Rossi et al.'s
improvisations and GuitarSet~\cite{xi2018guitarset}.

\begin{inapp}[: data used so far]
\begin{itemize}
\item \textbf{Benchmark}: EGDB test items (DI, the Marshall, Mesa and Plexi
  renders, and labelled notes), through the manifest \code{eval_egdb.tsv}
  written by \code{scripts/train/prepare_data.py} (Section~\ref{sec:results}).
  The working tree adds Guitar-TECHS miked-amp recordings and separated stems as
  targets \wip.
\item \textbf{Fitting the plucked-string model}: 14 Guitar-TECHS DI recordings
  (P1/P2 single notes, chords, scales), CC BY 4.0. Never EGDB, which is the
  benchmark.
\item \textbf{Sampler dry signal}: Guitar-TECHS ``allsinglenotes'' DI \wip.
\end{itemize}
All of it lives outside the repository, under \file{~/.cache/gva_test}.
\end{inapp}

% !TeX root = ../main.tex
\section{Effect recognition and parameter estimation}
\label{sec:recognition}

This is the most direct line of work: given processed audio, say which effects
are present and with which settings. Almost all of it trains on synthetic
renders, where the true chain is known.

\subsection{Single effects on isolated notes}

Stein et al.\ detected single and cascaded effects on IDMT notes with
hand-crafted spectral, cepstral and harmonic features and support vector
machines, reaching 99.2\% for combinations of six effects on single guitar notes
\reported~\cite{stein2010aes,stein2010dafx}. Such features transfer poorly:
from single notes to polyphonic material, recall fell to chance with one feature
set and to 41--62\% with a bag-of-audio-words model~\cite{schmitt2017recognising}.
The same study noted that delay is trivial to detect on single notes and hard on
phrases, where echoes look like repeated notes. Jürgens et al.\ added parameter
regression with shallow networks (MAE 0.05--0.16 on a unit scale for distortion,
delay and tremolo controls)~\cite{jurgens2020recognizing}. Comunità et al.\
trained a CNN (FxNet) to identify 13 drive plugins (81--91\% accuracy) and a
regressor for gain and tone (MAE 0.02--0.07)~\cite{comunita2021guitar}; most
confusions were between nearly identical circuits such as the TS808 and TS9, a
first sign that the inverse problem is ill-posed at the level of device
identity.

\subsection{Chains and mixes}

Hinrichs et al.\ estimated effect classes and parameters with the guitar still
inside a mix of bass, keys and drums: 97.4\% accuracy on isolated notes and
90.0\% in mixes \reported~\cite{hinrichs2022convolutional}. Guo and McFee
recognised combinations of 13 SoX effects on guitar phrases, with micro and macro
F1 of 0.876 and 0.906 on an unseen dataset~\cite{guo2023automatic}; overdrive vs
distortion, slapback vs feedback delay, and the modulation effects were the usual
confusions. \textbf{Order is harder than presence}: Wada et al.\ reach a set F1
of 0.987 on three-effect chains but a macro F1 of 0.756 on their 16
orders~\cite{wada2025hyperbolic}; Rossi et al.\ an exact-chain accuracy of 0.815
when order and parameters vary~\cite{rossi2025automatic}. Take et al.\ estimated
the chain while recovering the dry signal by peeling effects one at a time;
re-applying the estimated chain to the recovered dry signal improved the wet
reconstruction by only 0.67\dB{} SI-SDR, against 2.24\dB{} from the true dry
signal, because removal and estimation errors add up~\cite{take2024audio}.

\begin{keyidea}[: prediction plus search wins]
Okita and Katayose~\cite{okita2026audio} combined the two approaches of
Section~\ref{sec:search}: a network predicts the dry signal and the effect types,
then CMA-ES or TPE searches the order and parameters of a black-box chain to
maximise the SI-SDR of the re-rendered wet signal. With chorus, distortion and
reverb, the search raised the wet reconstruction from 18.18 to 23.07\dB{} and
lowered MR-STFT from 0.465 to 0.340 \reported. The chorus rate had to be fixed:
its landscape was too rugged. This is the closest published relative of the
app's estimator.
\end{keyidea}

\begin{table}[h]
\centering
\caption{Effect recognition and parameter estimation on guitar, as published.
All training data are synthetic renders.}
\label{tab:recognition}
\small
\begin{tabular}{@{}llll@{}}
\toprule
Work & Input & Task & Result \reported \\
\midrule
Stein~\cite{stein2010dafx} & notes & 6 cascaded effects & 99.2\% acc. \\
Comunità~\cite{comunita2021guitar} & notes & 13 drives, gain, tone & 81--91\%, MAE 0.04 \\
Hinrichs~\cite{hinrichs2022convolutional} & notes in mix & class, params & 90.0\%, MAE $<$0.016 \\
Guo~\cite{guo2023automatic} & phrases & 13 effects & F1 0.876 / 0.906 \\
Rossi~\cite{rossi2025automatic} & DI phrases & 5-effect chain & exact 0.815 \\
Wada~\cite{wada2025hyperbolic} & phrases & ordered chain & set 0.987, order 0.756 \\
Okita~\cite{okita2026audio} & phrases & type, then search & 23.07\dB{} SI-SDR \\
StemFX~\cite{cheng2026stemfx} & separated stems & mixing FX chain & top-1 86.8\% \\
\bottomrule
\end{tabular}
\end{table}

\subsection{Whole processing graphs}

Outside guitar, Lee et al.\ estimate a whole processing graph (nodes,
connections, parameters) from a processed reference~\cite{lee2023blind}; most
errors were wrong orders or processors hidden by a later destructive one. Later
they fit an over-complete differentiable console to dry tracks and the mix and
prune two-thirds of the processors without losing quality~\cite{lee2024searching,lee2025reverse},
using the GRAFX library~\cite{lee2024grafx}. Colonel and Reiss recovered gain,
pan, EQ, delay and reverb of a mix from its raw tracks~\cite{colonel2021reverse}.
These need the dry tracks, but one idea transfers directly: \emph{fit a fixed,
over-complete chain and let the fit switch parts off}. The app's template chain
does exactly that (a pedal whose amount falls below 0.04 is removed).
StemFX is the first system trained to predict effect chains directly on
separated stems without a dry reference: a band-split encoder and a transformer
decoder over 85 effects, trained on pseudo-stems of about 105k
songs~\cite{cheng2026stemfx}. It predicts mixing effects, not guitar rigs, and
guitar sits in its ``other'' stem. LLM2Fx-Tools predicts executable effect chains
as tool calls from dry and processed pairs~\cite{doh2026llm2fxtools}.

\subsection{Sub-problems of a guitar chain}

\begin{description}[style=nextline,leftmargin=1.2cm]
\item[Modulation] Mitcheltree et al.\ extract the LFO of phasers, flangers and
  choruses from processed audio with 1.8--4.9\% error on guitar and 3.1\% on
  unseen real guitar~\cite{mitcheltree2023modulation}. Carson et al.\ fit
  differentiable flangers, choruses and phasers and need a low-frequency
  weighting of the loss to avoid local minima in delay
  time~\cite{carson2026gradient}.
\item[Delay and reverb] Delay time and mix were estimated only on synthetic
  notes~\cite{jurgens2020recognizing,hinrichs2022convolutional}. Blind estimation
  of artificial reverb parameters exists for speech~\cite{lee2022differentiable},
  and blind RT60 estimation is a room-acoustics topic~\cite{gamper2018blind}; we
  found nothing on the reverb of a guitar track inside a mix.
\item[Dynamics] Compressor parameters by hybrid neural regression with
  model-based inversion~\cite{sun2024neural} or gradient-free search in a
  perceptual feature space~\cite{sun2026blackbox}.
\item[The instrument] Pickup position and plucking point can be estimated from
  the spectrum, also through amplifier chains (median errors of 2--7\,mm for the
  pickup)~\cite{mohamad2017pickup,mohamad2017estimating}. The guitar is part of
  the tone that an effects app cannot change, only compensate with EQ.
\end{description}

% !TeX root = ../main.tex
\section{Amplifier and effect modelling}
\label{sec:modelling}

Modelling is the \emph{forward} problem: given the dry signal, produce the
processed one. It matters here for two reasons: a model can be the output of a
tone estimator (a capture loaded into a player), and a model with few,
interpretable parameters makes the inverse problem easier to search.

\subsection{Black-box models from paired data}

\begin{keyidea}[: capture = learn the device from dry/wet pairs]
Record a few minutes of dry guitar, play it through the real amplifier, record
the output, and train a network to map one to the other with an ESR loss. The
network then imitates that amplifier at that setting. This is what Neural Amp
Modeler, Kemper profiling, TONEX and Neural Capture do.
\end{keyidea}

Feed-forward WaveNets~\cite{damskagg2019deep,damskagg2019realtime} and recurrent
networks~\cite{wright2019realtime} model amplifiers and pedals from a few minutes
of paired audio. Wright et al.\ report test ESR of 0.2--0.3\% and MUSHRA scores
of 90--98 on two tube amplifiers, faster than real time on a CPU
\reported~\cite{wright2020realtime}; a perceptual pre-emphasis improves the
result at no cost~\cite{wright2020perceptual}. Temporal convolutional networks
model compressors~\cite{steinmetz2022efficient}; time-varying feature modulation
extends their memory for fuzz and compression~\cite{comunita2023modelling}. The
largest comparison (40 devices; LSTM, TCN, GCN and S4 models) ranks state-space
models with time-varying modulation best and finds LSTMs unreliable across
devices~\cite{comunita2025differentiable}. Vanhatalo et al.\ review the
field~\cite{vanhatalo2022review}.

Neural Amp Modeler (NAM) has become the open format for captures: an MIT-licensed
trainer, core and plugin~\cite{atkinsonndnam}, a JSON \file{.nam} file with
WaveNet (``A1'', and the slimmable ``A2''~\cite{atkinson2025slimmable,atkinson2026a2})
and LSTM architectures, and optional metadata (gear type, tone type, gain). The
metadata are ready-made weak labels for a capture library. Any capture is one
setting of one device, made with the device present.

\begin{inapp}[: NAM in the app]
The NAM Model pedal runs \file{.nam} files through the pure-Rust \code{nam-rs}
runtime; a standard WaveNet runs at about 11 times real time on one x86 core
and does not resample, so models must match the 48\kHz{} host rate. The tone
estimator does not use NAM captures yet (Section~\ref{sec:design}).
\end{inapp}

\subsection{Knobs and many devices in one model}

A model can take the knobs as inputs. Juvela et al.\ recorded 4.5\,h of an
amplifier with a robot turning its knobs and trained a conditioned LSTM that
listeners could not tell from a circuit simulation~\cite{juvela2023endtoend}.
PANAMA chooses which settings to record by the disagreement of an ensemble, and
with 75 settings matches per-setting NAM captures in a listening
test~\cite{grotschla2025parametric}. Hypernetwork conditioning beats
concatenation for recurrent models~\cite{yeh2024hyper}.

One model can also cover many devices. Chen et al.\ condition a GCN on a tone
embedding from a contrastively trained encoder, so a wet reference clip selects
the tone; on unseen amplifiers the loss is about twice the seen-amplifier loss,
and a case study missed an EQ the model did not
represent~\cite{chen2024zeroshot,chen2024demo}. Open-Amp trains a TCN with a
394-entry device embedding on renders through crowd-sourced captures; a new
analog pedal is added by fitting only its embedding~\cite{wright2025openamp}.
Embedding-only enrolment is a small search space, a candidate target for a
wet-only loss (the paper used paired data).

\subsection{Grey-box and differentiable signal processing}

\term{Grey-box} models keep the structure of the circuit and learn only its
values. The classic block structure is \term{Wiener--Hammerstein}: a linear
filter, a static (memoryless) nonlinearity, another linear
filter~\cite{eichas2017block,eichas2018graybox}, later made
differentiable~\cite{kuznetsov2020differentiable,colonel2022waveshapers}.
Esqueda et al.\ learn the component values of Fender, Marshall and Vox tone
stacks from audio~\cite{esqueda2021differentiable}. Miklánek et al.\ put a
white-box tone stack between a recurrent preamp and power stage; trained on
4--12 minutes covering one to three tone-stack settings, it scored 87--95 MUSHRA
on unseen settings where a black-box network trained on the same data scored
18--47 \reported~\cite{miklanek2023neural}. The DDSP guitar amplifier of Yeh et
al.\ (preamp, tone stack, push--pull power amp, output transformer) needs about
1\,350 operations per sample and beats a recurrent model of the same
cost~\cite{yeh2025ddsp}. Grey-box compressors~\cite{wright2022greybox},
phasers~\cite{carson2023differentiable} and modulation
effects~\cite{carson2026gradient} follow the same idea. Grey-box models are not
yet as accurate as neural networks for distortion and fuzz, though competitive
for compressors~\cite{comunita2025differentiable}.

Libraries make these models reusable: DDSP~\cite{engel2020ddsp}, dasp-pytorch
and DeepAFx-ST~\cite{steinmetz2022style}, GRAFX~\cite{lee2024grafx},
NablAFx~\cite{comunita2025nablafx}, DiffVox~\cite{yu2025diffvox}; Hayes et al.\
review differentiable signal processing~\cite{hayes2024review}. DiffVox fits a
differentiable EQ, compressor, delay and reverb to each of 435 vocal tracks and
studies the distribution of the fitted presets, a recipe that would turn guitar
stems into a prior over settings.

\begin{keyidea}[: why this matters for estimation]
The app's amplifier is itself grey-box-like: gain, a voicing, a three-band tone
stack. That matters less for modelling quality than for estimation: a few
interpretable parameters are what a search can explore and what a player can
adjust afterwards.
\end{keyidea}

% !TeX root = ../main.tex
\section{Working without the dry signal}
\label{sec:nodi}

Everything in Section~\ref{sec:modelling} needs the dry signal of the same
performance. A song does not have it. There are three ways around this.

\subsection{Unpaired and blind modelling}

\begin{keyidea}[: learn the device from its output's statistics]
If we cannot compare the model's output with the target sample by sample, we
can still ask that the model's output \emph{looks like} processed guitar in
general. A \term{discriminator} network trained to tell real processed guitar
from the model's output provides that signal (adversarial training, as in GANs).
\end{keyidea}

Wright, Välimäki and Juvela trained a real-time WaveNet amp model only
adversarially~\cite{wright2023adversarial}. With input and target from the same
guitar, the best unpaired model scored 83--92 MUSHRA across three plugin tones;
with a different guitar for the input, the realistic case, 77--84, comparable to
an oracle that equalised the true plugin's output (72--81) \reported. Training
took 4--10 GPU hours per tone. Chen et al.\ improved stability with multi-scale
and multi-period discriminators and extra clean data, and found loudness
normalisation essential~\cite{chen2024unsupervised}.

Moliner et al.\ estimate the clean signals and the effect jointly by
\term{diffusion posterior sampling} with a prior trained on clean
guitar~\cite{moliner2025unsupervised}. Diffusion stays stable with only 18\,s of
processed audio, where adversarial training degrades; adversarial training is
better on heavy distortion; a grey-box operator matched black-box ones under
diffusion but collapsed under adversarial training. Švento et al.\ estimate a
memoryless nonlinearity (a spline) together with the clean
signal~\cite{svento2025estimation}, giving both a dry estimate and a waveshaper
curve; LOUDAR moves the operator into an autoencoder's latent
space~\cite{svento2026loudar}.

All of these were evaluated on isolated, plugin-rendered guitar, mostly with
distortion alone. None handles cabinet, modulation, delay and reverb together,
and none was tested on a separated stem.

\subsection{Effect removal}

A recovered dry signal turns the wet-only problem back into a paired one. Imort
et al.\ removed distortion with separation networks~\cite{imort2022distortion}.
RemFX detects which of five effects are present and chains effect-specific
removers; SI-SDR falls from 16.7\dB{} with one effect to 2.6\dB{} with
five~\cite{rice2023general}. Distortion Recovery restores a mel spectrogram and
re-synthesises it with a vocoder~\cite{lee2024distortion}. EG-VAE splits a
processed recording into frame-level content and a global tone embedding;
swapping the embedding transfers the tone, a masking mode removes
it~\cite{yeh2026egvae}. On plugin tones it reaches a tone-similarity opinion
score of 4.30 on seen and 3.70 on unseen tones (oracle 4.48; baselines at most
2.74), and a ``dryness'' score of 3.42--3.50 for removal against 4.39--4.58 for
the true DI \reported. A baseline that fitted a fixed differentiable chain of
generic blocks to EG-VAE's dry estimate scored only 1.44--2.74 in tone
similarity: evidence that generic EQ, compressor and distortion blocks do not
cover guitar tone diversity. No code or weights were found for EG-VAE.

Removal and remodelling have been chained: Hinrichs et al.\ remove a single
effect with a separation network, then train a model on the recovered pair;
listeners rated the copies at medians of 92--100 for slapback, phaser and soft
clipping~\cite{hinrichs2024blind}. SDR does not track perceived dryness for
distortion, as also found by Imort et al.\ and the EG-VAE authors.

\subsection{Dry signals from other sources}

Three other sources of a dry signal exist, each with its own mismatch:
\begin{center}
\small
\begin{tabular}{@{}>{\raggedright}p{4.2cm}ccc@{}}
\toprule
Source & Same notes and timing & Same guitar timbre & Artefacts \\
\midrule
A different performance (corpus DI, the player's DI) & no & roughly & none \\
Resynthesis of the transcribed notes & yes & no & transcription errors \\
An effect-removal model & yes & yes & removal errors \\
\bottomrule
\end{tabular}
\end{center}
A different performance supports only unaligned losses, as in unpaired
modelling~\cite{wright2023adversarial}. Resynthesis gives the same notes at the
same times with the wrong timbre; GuitarFlow turns sample-based renders of
tablature into realistic DI with flow matching~\cite{loth2025guitarflow}, but for
data generation, not tone estimation. \textbf{We found no published work that
resynthesises a DI from transcribed notes in order to fit a tone.} That is the
route the app takes first (Section~\ref{sec:impl-dry}), because the app already
transcribes the notes.

% !TeX root = ../main.tex
\section{Style transfer, effect representations and sound matching}
\label{sec:matching}

\subsection{Reference-based style transfer}

Audio-effect \term{style transfer} processes an input so that it sounds like a
reference. DeepAFx-ST predicts EQ and compressor parameters from an input and a
reference, trained self-supervised with an audio loss; numerical gradients (SPSA)
let it drive non-differentiable effects at a cost in
stability~\cite{steinmetz2022style}. Grant trained the same way through legacy
plugins~\cite{grant2023style}. Koo et al.\ learn an effects encoder by
contrastive learning and render input stems in a reference's mixing style; its
embedding distances on Demucs-separated stems stayed close to those on clean
stems~\cite{koo2023music}. Diff-MST predicts a differentiable mixing console for
any number of tracks~\cite{vanka2024diffmst}.

\textbf{ST-ITO} removes the trained controller: CMA-ES searches the parameters of
an arbitrary plugin chain to maximise the cosine similarity of AFx-Rep embeddings
of the processed input and the reference~\cite{steinmetz2024stito}. AFx-Rep is
trained without dry references to recognise which of 63 plugins and 10 presets
produced an output, and beats CLAP and MFCCs on zero-shot style classification
(0.86 against 0.56 for CLAP) \reported. The search uses a population of 64 for at
most 25 steps (about a minute). The authors state that it does not work well for
challenging cases such as \emph{guitar tone matching}. Yu et al.\ add a Gaussian
prior fitted on real presets, turning the search into a maximum-a-posteriori
estimate and reducing parameter error by up to a third~\cite{yu2025improving}.
ITO-Master found that optimising a learned conditioning embedding beats
optimising raw DSP parameters for mastering~\cite{koo2025itomaster}.

\begin{caution}[: an embedding is not an objective]
TAPE predicts editable parameters of its own differentiable chain (EQ,
compressor, waveshaper, a small residual network, reverb) in one pass and runs in
real time~\cite{long2026tape}. Its comparison is the most instructive result for
us: CMA-ES over ST-ITO's chain with the AFx-Rep objective lowered the embedding
distance (0.259 against 0.660 for TAPE) but raised the spectral error to 2.91,
\emph{above} the 2.34 of the unprocessed input, while TAPE reached 1.12
\reported. An embedding that is good for retrieval can be a bad objective to
optimise. The app's search therefore uses a spectral distance.
\end{caution}

\subsection{Effect representations}

Several encoders describe effects independently of content: AFx-Rep, the effects
encoder of Koo et al., Fx-Encoder++, the tone encoder of Chen et
al.~\cite{chen2024zeroshot}, the capture-trained encoder of Open-Amp (91.1\%
effect-class accuracy on EGDB~\cite{wright2025openamp}) and TAPE's.
Fx-Encoder++ extracts an instrument-wise effects embedding from a mixture given
an audio or text query~\cite{yeh2025fxencoder}, and gives \textbf{the only
measurement on separated stems}: effect retrieval R@10 falls from 39.4\% on true
isolated stems to 20.5\% on Demucs-separated stems \reported, attributed to
separation artefacts distorting effect cues. RelFx learns the \emph{relative}
effect between two processed clips, so it needs no dry
reference~\cite{liu2026beyond}. General embeddings (OpenL3, PANNs, CLAP) move
monotonically with effect strength but not linearly, so they are weak parameter
estimators~\cite{deng2025investigating}. TimberAgent retrieves the nearest of
1\,063 guitar-effect presets with Gram matrices of Wav2Vec2
features~\cite{he2026timberagent}.

\subsection{Text and language models}

Text2FX optimises EQ and reverb against a CLAP text embedding~\cite{chu2025text2fx};
LLM2Fx predicts them with language models, better than Text2FX on guitar EQ
(distribution distance 0.15 against 0.42)~\cite{doh2025llm2fx}; FxSearcher uses
Bayesian optimisation over a plugin chain~\cite{ki2025fxsearcher}. The ``Does it
chug?'' dataset collects guitar-tone adjectives on clips that vary EQ and
distortion~\cite{sutar2024chug}. None estimates amp or drive settings; text suits
refinement (``more bite'') better than estimation.

\subsection{Lessons from synthesizer sound matching}

Finding synthesizer settings that reproduce a sound is the same inverse problem
with a longer history. Its findings carry over:
\begin{itemize}
\item \textbf{Loss choice}: training on audio losses alone often fails for
  frequency-like parameters; pre-training on a parameter loss and fine-tuning on
  a spectral loss works best~\cite{masuda2021ddsp,uzrad2024diffmoog}. LFO rates,
  delay times and wah positions are such parameters.
\item \textbf{Optimisers}: on non-convex landscapes CMA-ES beat Adam (loss 2.09
  against 4.60)~\cite{bogdan2026instrumental}; evolution strategies match sounds
  through non-differentiable synthesizers with a CLAP objective~\cite{cherep2024ctag};
  reinforcement learning with rewards from rendered audio improves
  out-of-domain matching~\cite{wu2026ddsynthrl}.
\item \textbf{Ambiguity}: when different settings give the same sound, point
  regression fails and generative models of the parameters do
  better~\cite{hayes2025equivariant}.
\item \textbf{Proxies}: a network that maps presets to a pretrained audio
  embedding makes black-box synthesizers trainable~\cite{combes2025proxies}.
\item \textbf{Domain gap}: every system degrades on real sounds outside its
  training renders~\cite{hayes2024review}.
\end{itemize}

\begin{keyidea}[: the consensus]
The effect and synthesizer literatures agree on a hybrid: predict an initial
setting with a network trained on renders of the target chain, refine it by
gradient-free search through the real renderer, regularise with a prior over
realistic presets, and keep several candidates.
\end{keyidea}

% !TeX root = ../main.tex
\section{Products and open-source tools}
\label{sec:products}

Statements in this section come from manuals and product pages; vendors publish
no accuracy figures for recording-based tone features. Products take three kinds
of input (Table~\ref{tab:products}).

\paragraph{Capture of physical equipment.} Kemper profiling plays test signals
into the amplifier and analyses its response; compressors and time-based effects
must be bypassed~\cite{kemper2018manual}. Neural DSP's Neural Capture, IK
Multimedia's TONEX (with parametric models announced in October 2026), Line 6's
Proxy and NAM work the same way with neural
models~\cite{neuraldsp2026qcmanual,ik2026parametric,line6ndproxy,atkinsonndnam}.
None accepts a recording. The Quad Cortex manual says that a capture's gain
metadata is estimated from the device's dynamic range and harmonic content, the
only vendor statement found of a gain descriptor computed from
audio~\cite{neuraldsp2026qcmanual}.

\paragraph{Spectral matching to a recording.} The Axe-Fx III \textbf{Tone Match}
block computes long-term average spectra of a reference and of the user's preset
and applies their smoothed difference as a filter, exportable as a cabinet
IR~\cite{fractalndtonematch}. Its manual lists the conditions: a starting preset
with the right amplifier and the right amount of distortion, no effects, an
isolated and unclipped guitar track, matched loudness, mono analysis, broad
playing material. iZotope's Match EQ does the same for mixes~\cite{izotope2019matcheq}.
These tools estimate only the linear part, after the user has fixed the
nonlinear one. The app's final step is the same operation
(Section~\ref{sec:impl-match}), applied after its own search has fixed the gain.

\paragraph{Presets from audio, song names or text.} Positive Grid's BIAS X
``Music-to-Tone'' accepts a file of up to 15 minutes, full mixes included, finds a
guitar segment and returns a full chain; it runs online and no method or accuracy
is published~\cite{positivegridndbiasxhelp,positivegrid2025biasxai}. Small apps
generate presets from a song name with a language model's knowledge of the
artist's gear, sometimes with stem separation and coarse spectral
features~\cite{bouharb2026gptonebuilder,aitonetechndgenerator}; none publishes an
evaluation. The most detailed public design is a Chinese patent granted in 2025:
separation, a tone-type classifier, a network that picks an amp and cabinet model
and regresses five knobs, a fit of up to 24 peaking filters to the difference of
the spectra, tempo and key analysis for delay and modulation times, and retrieval
from a preset database for the remaining effects~\cite{huanyin2025presetpatent}.

\paragraph{Open-source tools.} NAM-EQ-Match renders the player's DI through every
capture in a NAM library at several input gains, compares EQ-invariant dynamics
features (crest factor, compression, spectral flux, MFCC statistics) with the
target, then fits a matching IR or tone-stack EQ~\cite{sd4562026nameqmatch}.
\code{tonematch} predicts seven amp knobs with a mel-spectrogram CNN and refines
them by CMA-ES through a grey-box amp, reporting a median log-spectral error of
1.54\dB{} on 24 isolated Guitar-TECHS cases against 3.68\dB{} with every knob at
noon~\cite{neuwirth2026tonematch}. ToneRunner retrieves plugin chains with a
contrastive encoder and documents how the encoder learns the riff instead of the
tone when the same dry clip is not shared across chains in
training~\cite{nathanben2026tonerunner}. These projects are small and not peer
reviewed, but they converge on the same design as the research: retrieve or
predict a nonlinear core, refine by rendering, finish with a linear match.

\begin{table}[h]
\centering
\caption{What products and tools take as input and return. ``Full chain'':
amplifier, cabinet and effects with settings. $^*$Method undocumented.}
\label{tab:products}
\small
\begin{tabular}{@{}lll@{}}
\toprule
System & Input & Output \\
\midrule
Capture products, NAM & the device itself & capture \\
Tone Match~\cite{fractalndtonematch} & isolated track & EQ / IR \\
BIAS X~\cite{positivegridndbiasxhelp} & recording, even a mix & full chain$^*$ \\
Song-name apps~\cite{bouharb2026gptonebuilder} & name, text & preset \\
CN118865922B~\cite{huanyin2025presetpatent} & recording & full chain \\
NAM-EQ-Match~\cite{sd4562026nameqmatch} & stem + player's DI & capture, EQ/IR \\
\code{tonematch}~\cite{neuwirth2026tonematch} & recording + DI & 7 amp knobs \\
Guitar Virtual Amp (this work) & separated stem + notes & app preset + match IR \\
\bottomrule
\end{tabular}
\end{table}

% !TeX root = ../main.tex
\section{Comparison, failure modes and open problems}
\label{sec:comparison}

\subsection{Comparison}

The commercial case needs three properties at once: no dry signal of the target,
an output a real-time chain can use, and robustness to separated stems.
Table~\ref{tab:methods} shows that no published method has all three. Effect
recognition returns usable parameters but was trained and tested almost only on
synthetic, isolated renders. Unpaired modelling needs no aligned dry signal and
returns a real-time model, but covered only distortion on isolated guitar. Tone
transfer (EG-VAE) handles whole chains but returns audio. Inference-time search
(ST-ITO) drives any chain but fails on guitar tone with an embedding objective.
StemFX is the only work trained on separated stems, for mixing effects.

\begin{table}[h]
\centering
\caption{Main methods by requirement and output. ``Non-diff.'': works with a
chain that cannot be differentiated. ``Sep.'': trained or tested on separated
stems. Licence of released code or weights (NC = non-commercial; -- = not
released or no licence file; ? = code released, licence not checked).}
\label{tab:methods}
\footnotesize
\setlength{\tabcolsep}{2.5pt}
\resizebox{\textwidth}{!}{%
\begin{tabular}{@{}llllllll@{}}
\toprule
Method & Needs & Output & Non-diff. & Guitar & Sep. & Code & Reported \\
\midrule
FxNet~\cite{comunita2021guitar} & wet & drive type, gain, tone & yes & yes & no & BSD-3 & 81--91\%, MAE $\approx$0.04 \\
Hinrichs 2022~\cite{hinrichs2022convolutional} & wet in mix & type, params & yes & yes & no (mix) & -- & 90.0\% in mixes \\
Okita~\cite{okita2026audio} & wet & chain (search) & yes & yes & no & -- & 23.07\dB{} wet SI-SDR \\
StemFX~\cite{cheng2026stemfx} & wet stems & mixing FX chain & yes & no & yes & MIT & 86.8\% top-1 \\
DeepAFx-ST~\cite{steinmetz2022style} & wet + input & EQ, comp. & SPSA & no & no & NC & -- \\
ST-ITO~\cite{steinmetz2024stito} & wet + input & any chain (search) & yes & fails & no & Apache-2.0 & 0.86 style acc. \\
TAPE~\cite{long2026tape} & wet + input & own DSP chain & no & partly & no & -- & STFT L1 1.12 vs 2.34 \\
Fx-Encoder++~\cite{yeh2025fxencoder} & mix + query & embedding & -- & partly & yes & NC & R@10 39.4 $\to$ 20.5\% \\
Unpaired GAN~\cite{wright2023adversarial} & wet + other DI & neural amp & -- & yes & no & -- & MUSHRA 77--84 \\
Diffusion blind~\cite{moliner2025unsupervised} & wet + clean prior & neural / grey-box & -- & yes & no & ? & stable from 18\,s \\
Remove + model~\cite{hinrichs2024blind} & wet & neural effect & -- & yes & no & -- & MUSHRA $\geq$92 (1 effect) \\
EG-VAE~\cite{yeh2026egvae} & wet + any DI & audio, DI estimate & -- & yes & no & -- & MOS 4.30 / 3.70 \\
Tone embedding~\cite{chen2024zeroshot} & wet ref.\ + DI & neural amp & -- & yes & no & -- & 2$\times$ loss unseen \\
Open-Amp~\cite{wright2025openamp} & paired (enrol) & neural, 394 devices & -- & yes & no & -- & 91.1\% class acc. \\
Tone Match~\cite{fractalndtonematch} & isolated wet & EQ / IR & yes & yes & no & product & -- \\
\midrule
This app & stem + notes & app chain + IR & yes & yes & \wip & GPL-3.0 & Section~\ref{sec:results} \\
\bottomrule
\end{tabular}}
\end{table}

\subsection{Known and expected failure modes}

\begin{description}[style=nextline,leftmargin=1.2cm]
\item[Separation artefacts] Effect retrieval roughly halves from isolated to
  separated stems \reported~\cite{yeh2025fxencoder}. Reverb tails, stereo delays
  and wide choruses are the most exposed to leakage and masking; bleed from keys
  or other guitars adds energy a matcher will try to reproduce \expected.
\item[Several guitars in one stem] Double-tracked or layered guitars give one
  stem with mixed tones. The pan split separates parts on different sides, not
  in the same position \expected.
\item[Mixing and mastering] Bus compression, limiting and mix EQ act after the
  guitarist's chain. An estimator attributes them to the guitar; a final linear
  match absorbs the EQ but not the dynamics \expected.
\item[Non-uniqueness] Drive vs amp gain, fuzz vs high-gain amp, compression vs
  picking dynamics, effect order. Device identity is confused even between clean
  renders of similar circuits~\cite{comunita2021guitar}, and order is harder than
  presence~\cite{wada2025hyperbolic}. Harmless for reproduction if the result
  sounds alike, which is why audio objectives are preferred~\cite{peladeau2024blind}.
  The app's benchmark shows it: the right amp voicing is found in only 7 of 12
  items, yet the sound gets closer \measured.
\item[Equipment outside the chain] Guitar, pickups and microphone shape the tone
  and can only be compensated by EQ. A mono chain cannot reproduce stereo
  effects. Generic blocks may lack a specific amp's character, as EG-VAE's
  differentiable baseline showed~\cite{yeh2026egvae}.
\item[Search pathologies] LFO rates and delay times make rugged
  landscapes~\cite{okita2026audio,uzrad2024diffmoog}; an embedding objective can
  improve while the sound moves away~\cite{long2026tape}; errors of a removal
  model add to errors of the estimate~\cite{take2024audio}.
\item[Tone changes within a song] Clean verses and distorted choruses need
  different presets. Segmenting songs into tone sections exists only in
  products~\cite{bouharb2026gptonebuilder}.
\end{description}

\subsection{Open problems}

\begin{enumerate}
\item Estimating an interpretable amp, cabinet and effects chain from a separated
  stem of a commercial song, evaluated against known equipment.
\item A benchmark for that task: commercial-like mixes with known chains, and
  real multitracks with the guitarist's DI and amp tracks, scored with metrics
  validated by listening tests.
\item Unpaired or blind modelling of a full chain, on separated rather than
  isolated guitar.
\item Delay, reverb and modulation estimation on real guitar phrases, with
  validated metrics.
\item Training data that can be redistributed: most realistic sets and the
  largest capture libraries are non-commercial or forbid bulk use.
\end{enumerate}


\part{Guitar Virtual Amp}
% !TeX root = ../main.tex
\section{The design for Guitar Virtual Amp}
\label{sec:design}

This section is the plan derived from Part~II, written before any code. Each
step carries its status today; Section~\ref{sec:impl} describes what is built
in detail.

\subsection{What the app already had}

\begin{itemize}
\item \textbf{A guitar stem per part}: BS-RoFormer (Demucs as fallback)
  separation~\cite{lu2023bsroformer,rouard2023htdemucs} and a split of the stereo
  guitar stem by pan position into parts that play different
  things~\cite{deangelo2026tabreview}.
\item \textbf{The notes and the grid}: onsets, offsets, pitches and techniques of
  every guitar part; beats, bars and tempo from Beat This!~\cite{foscarin2024beatthis}.
\item \textbf{An offline renderer of the real chain}: \code{backend::offline::render_at}
  runs the same pedal code as the live path (48\kHz, 16-bit samples, 256-sample
  blocks) at a chosen input RMS level. The full chain of 17 classic pedals costs
  about 430\,ns per sample on one core, a NAM WaveNet about 1.9\,\textmu s.
\item \textbf{Presets with song metadata}: 20 factory presets, user presets with
  \code{song} and \code{artist} fields, and a tone suggestion on the Tabs page
  that prefers a user preset tagged with the song (otherwise it guesses a factory
  preset from the track's instrument and name).
\item \textbf{Capture and cabinet players}: the NAM pedal and an IR convolver.
\end{itemize}

\subsection{Constraints}

\begin{itemize}
\item The chain is mono, works on 16-bit samples and has no gradients: the
  estimator treats it as a black box. Stereo effects cannot be reproduced.
\item The estimate may run offline on a desktop; its output, a preset, also
  works on Android.
\item The analysed audio must be the user's own: Spotify's developer policy
  forbids analysing its content, so the feature applies to the same local files
  and downloads as transcription.
\item Models that are distributed must be trained on data that allow it
  (Section~\ref{sec:datasets}); capture libraries can be searched by metadata,
  not bulk downloaded.
\item Patent CN118865922B claims a pipeline of similar shape~\cite{huanyin2025presetpatent};
  its claims should be reviewed before distributing the feature where the patent
  applies.
\end{itemize}

\subsection{The planned pipeline}

\begin{figure}[h]
\centering
\begin{tikzpicture}[
  node distance=2.4mm,
  box/.style={draw, rounded corners=2pt, font=\small, align=left,
              minimum height=7mm, inner sep=3pt, fill=white, text width=125mm},
  arr/.style={-{Stealth[length=2mm]}, semithick}]
  \node[box, fill=appc!10] (s1) {1. Separation, pan split, notes, beats \hfill \textsf{\scriptsize built (transcription)}};
  \node[box, below=of s1, fill=appc!10] (s2) {2. Pick an excerpt per part \hfill \textsf{\scriptsize built: loudest 12\,s}};
  \node[box, below=of s2, fill=appc!10] (s3) {3. Stand-in dry signal from the notes \hfill \textsf{\scriptsize built: pluck, calibrated on the player's guitar}};
  \node[box, below=of s3, fill=warnc!10] (s4) {4. Starting points: factory presets; learned guess \hfill \textsf{\scriptsize presets built; learned planned}};
  \node[box, below=of s4, fill=appc!10] (s5) {5. CMA-ES through the offline renderer, preset prior \hfill \textsf{\scriptsize built}};
  \node[box, below=of s5, fill=warnc!10] (s6) {6. Delay, reverb, modulation from notes and tempo \hfill \textsf{\scriptsize delay built}};
  \node[box, below=of s6, fill=appc!10] (s7) {7. Linear match: minimum-phase IR in a Cabinet \hfill \textsf{\scriptsize built}};
  \node[box, below=of s7, fill=appc!10] (s8) {8. Save as presets tagged with the song and track \hfill \textsf{\scriptsize built}};
  \foreach \a/\b in {s1/s2,s2/s3,s3/s4,s4/s5,s5/s6,s6/s7,s7/s8} \draw[arr] (\a) -- (\b);
\end{tikzpicture}
\caption{The planned tone-estimation pipeline and the status of each step.
Green: built and committed; orange: partly built or in progress.}
\label{fig:pipeline}
\end{figure}

The plan, step by step (Figure~\ref{fig:pipeline}):
\begin{enumerate}
\item \textbf{Separation and transcription} come from the tablature pipeline.
\item \textbf{Excerpts}: for each part, pick 5--20\,s where the part is loud
  relative to the rest of the mix, other guitars are quiet and the notes cover
  several registers; cluster excerpts by an effect embedding to find tone changes
  between sections and estimate one preset per cluster. \emph{Built so far}: the
  loudest 12\,s window.
\item \textbf{Stand-in dry signal}, from (a) a resynthesis of the transcribed
  notes, (b) an effect-removal model trained on the app's own renders, or (c) the
  player's own DI recorded while playing along in practice mode. (c) has the right
  guitar and the wrong timing; (a) the right timing and the wrong guitar; (b)
  both, with removal errors. \emph{Built}: (a). A fourth source, measured
  later, is a mix: (a) for the notes, with its colour calibrated on (c)
  (Section~\ref{sec:results-colour}).
\item \textbf{Starting points}: the closest factory preset, then a network
  trained on renders of the template chain, mixed over backing tracks and
  separated by the same BS-RoFormer so it sees separation artefacts (the StemFX
  recipe~\cite{cheng2026stemfx}), sharing each dry clip across many chains so it
  learns the tone rather than the riff~\cite{nathanben2026tonerunner}. AFx-Rep
  (Apache-2.0)~\cite{steinmetz2024stito} is a starting point for its encoder.
\item \textbf{Refinement} by CMA-ES. The plan lists four terms for the
  objective: an aligned spectral term, an alignment-free term (long-term spectra
  per register, crest factor, compression statistics), a small-weight
  effect-embedding term, and a prior over factory and user
  presets~\cite{yu2025improving}. \emph{Built}: the aligned spectral term, with an
  EQ allowance, and a prior over the factory presets.
\item \textbf{Time-based effects from the notes}: delay time from echoes of
  onsets (or candidate values tied to the tempo), reverb decay from rests after a
  note ends, modulation rate from sustained notes or an LFO
  extractor~\cite{mitcheltree2023modulation}. Measured values fix or initialise
  the search, which avoids rugged landscapes. \emph{Built}: delay.
\item \textbf{Linear match}: the smoothed difference between the long-term
  spectra of the stem and of the rendered chain, as a minimum-phase IR in a
  Cabinet pedal with a gain limit, as Tone Match does~\cite{fractalndtonematch}.
\item \textbf{Output}: a user preset with song and artist set, so the existing
  tone suggestion offers it; other candidates kept; settings recorded for later
  evaluation.
\end{enumerate}
\textbf{Captures} (later): NAM files on the user's disk as candidates for the amp
slot, ranked by rendering the dry signal through each at several input
gains~\cite{sd4562026nameqmatch}; TONE3000 queried by metadata for candidates the
user then chooses to download. Learning a capture from the stem alone by unpaired
or diffusion-based fitting~\cite{wright2023adversarial,moliner2025unsupervised}
is research, at hours of GPU time per tone.

\subsection{Why this design}

\begin{itemize}
\item \textbf{Against prediction only}: synthetic-to-real gaps and ambiguous
  settings; refinement by rendering gained 4.9\dB{} SI-SDR in Okita and
  Katayose~\cite{okita2026audio}.
\item \textbf{Against an embedding-only search}: ST-ITO's authors report failure
  on guitar tone, and TAPE shows the objective drifting~\cite{steinmetz2024stito,long2026tape};
  a dry signal aligned with the target allows a spectral objective instead.
\item \textbf{Against neural tone transfer}: EG-VAE gives the best guitar tone
  similarity but returns audio, has no released code, and cannot run in the
  real-time chain~\cite{yeh2026egvae}.
\item \textbf{Against linear matching only}: a filter cannot set gain or
  dynamics, and Tone Match's own manual requires the gain to be right
  first~\cite{fractalndtonematch}.
\item \textbf{For the app's own chain}: its parameters are what the player
  edits, its renderer is the exact forward model, and training on its renders
  gives labels in its own parameter space with licences under our control.
\item \textbf{For a dry signal made from the notes}: the transcription already
  exists, and it gives the one thing the other stand-ins lack, alignment with the
  target. No published work does this.
\end{itemize}

% !TeX root = ../main.tex
\section{What is implemented}
\label{sec:impl}

All code is in \file{guitar_core/src/tone/}. It is blocking and CPU-heavy, so it
runs on worker threads only. This section follows the data through it, with the
constants the code uses, so that a change in the code can be checked against
this description.

\subsection{Overview}

\begin{inapp}[: entry points]
\begin{itemize}
\item \code{tone::estimate(target, dry, &Options)} \done: \code{target} and
  \code{dry} are 48\kHz{} mono and aligned. Returns an \code{Estimate}: the best
  search result (\code{best}), the closest factory preset (\code{preset}), the
  template settings, the excerpt used, the measured echo (\wip) and the number
  of renders.
\item \code{tone::estimate_from_notes(target, notes, &Options)} \done: renders the
  stand-in dry signal from notes with the plucked-string model, then calls
  \code{estimate}.
\item \code{tone::estimate_from_transcription(audio, notes, &Options)} \wip: the
  same, from the transcriber's notes (the model's mean posterior stands in for
  velocity: $v = 0.4 + 0.6\,a$, clamped to $[0.05, 1]$).
\item \code{tone::save_preset(estimate, title, artist, track)} \wip: writes the
  match IR and a user preset (Section~\ref{sec:impl-save}).
\end{itemize}
\code{Options} defaults: excerpt 12\,s, population 16, 30 generations, all CPU
threads, seed 1, spectral floor 80\dB.
\end{inapp}

The algorithm, in order:
\begin{enumerate}
\item Cut the 12\,s window of the target with the most energy (energy summed in
  0.1\,s blocks; \code{loudest_window}), and the same window of the dry signal.
\item Analyse the target once (\code{features::Target}).
\item Measure the delay time, if there is an echo (\code{analysis::echo_delay},
  \wip).
\item Score every factory preset at input RMS 0.1 (presets with a Looper or a
  Pitch Shifter are excluded); keep the closest as the baseline.
\item For each amp voicing (Clean, Crunch, Lead), run separable CMA-ES over the
  template chain from a neutral start; keep the best of the three.
\item Return the best chain and the EQ it still needs, which becomes the match
  IR.
\end{enumerate}

\subsection{The stand-in dry signal}
\label{sec:impl-dry}

\begin{keyidea}[: Karplus--Strong]
A plucked string can be simulated with a delay line whose length is one period
of the note, filled with a short noise burst and fed back through a gentle
low-pass filter: each pass around the loop the burst repeats (giving the pitch)
and loses some treble (giving the natural decay)~\cite{karplus1983digital}. The
\emph{extended} version~\cite{jaffe1983extensions} adds a fractional delay for
exact tuning, a decay-time control and a pluck-position comb filter.
\end{keyidea}

\begin{inapp}[: the plucked-string model, \file{tone/pluck.rs} \done]
For each note (start, end, MIDI pitch, velocity):
\begin{itemize}
\item \textbf{Delay line} of $f_s/f_0 - 0.5$ samples (the two-point loss filter
  adds half a sample), with linear interpolation for the fractional part.
\item \textbf{Excitation}: white noise through a one-pole low-pass (coefficient
  \code{brightness}, plus \code{velocity_brightness} times the velocity), minus a
  copy delayed by \code{pluck_position} times the period (the comb notches a pick
  at that point along the string creates), scaled by velocity$^{1.5}$.
\item \textbf{Loop}: two-point average (high-frequency loss) times a gain $g$
  chosen so the note decays by 60\dB{} in $T_{60}$:
  $g = 10^{-3/(f_0 T_{60})}$. While the note is held,
  $T_{60} = T_{\mathrm{lowE}}\cdot(f_{E2}/f_0)^{0.7}$ with $T_{\mathrm{lowE}}$ = \code{low_e_t60}, clamped to
  $[0.4\,\text{s}, T_{\mathrm{lowE}}]$ (higher notes ring shorter); after the note
  ends, $T_{60}$ = \code{release_secs}.
\item \textbf{Post}: a DC blocker, then the fixed filter \code{DI_EQ}: the
  measured long-term spectrum of real DI minus this model's, as a minimum-phase
  IR.
\end{itemize}
Settings were fitted by CMA-ES (\code{eval::fit_pluck}) on 14 Guitar-TECHS DI
recordings, minimising the EQ-forgiven distance to the real DI of the same
notes: \code{pluck_position} 0.48, \code{release_secs} 0.45, \code{low_e_t60} 16\,s,
\code{brightness} 0.008, \code{velocity_brightness} 0. The distance fell from
13.05 to 9.90\dB{} on the fitting set \measured. Refit only on Guitar-TECHS,
never on EGDB (the benchmark).
\end{inapp}

\begin{inapp}[: the sampler, \file{tone/sampler.rs} \wip]
An alternative stand-in built from real recordings: Guitar-TECHS
``allsinglenotes'' DI (every pitch of a guitar, plucked and held; CC BY 4.0).
Each transcribed note gets the recorded note of the nearest pitch, re-pitched by
resampling if needed, scaled by the velocity ratio (clamped to 0.2--4), started
10\ms{} early so the attack is whole, and faded out over 60\ms{} after the note's
end. It has a real pickup's attack and decay, from another guitar. Selected in
the benchmark with \code{GVA_TONE_SAMPLES}; not yet measured.
\end{inapp}

\subsection{The distance}
\label{sec:app-distance}

\file{tone/features.rs} \done. Both signals are first scaled to RMS 1.

\paragraph{Log-band spectrogram.} At each of two resolutions $r$ the signal is cut
into Hann-windowed frames, transformed, and the power is summed in triangular
bands spaced evenly on a log-frequency axis up to 16\kHz:
\begin{center}
\small
\begin{tabular}{@{}lccc@{}}
\toprule
Resolution & FFT size / hop (at 48\kHz) & Lowest band edge & Bands \\
\midrule
coarse time, fine frequency & 2048 / 512 (43\ms{} / 11\ms) & 60\Hz & 40 \\
fine time, coarse frequency & 512 / 128 (11\ms{} / 2.7\ms) & 250\Hz & 20 \\
\bottomrule
\end{tabular}
\end{center}
Band energies in dB are floored at 80\dB{} below the loudest band (the search
may use a higher floor, \code{Options::floor_db}, \wip). Call the target's value
$T_r(t,b)$ and the render's $R_r(t,b)$ for frame $t$ and band $b$.

\paragraph{Active frames.} Only frames of the target within 45\dB{} of its
loudest frame count (set $\mathcal{A}_r$); silence is ignored.

\paragraph{EQ allowance.} The long-term spectra are the means over active frames,
$\bar T_r(b)$ and $\bar R_r(b)$. Their difference
$\Delta_r(b) = \bar T_r(b) - \bar R_r(b)$ is split into its mean
$o_r = \overline{\Delta_r}$ (a level offset) and a centred, capped EQ
\[
  e_r(b) = \operatorname{clip}\bigl(\Delta_r(b) - o_r,\ -18, +18\bigr)\ \text{dB}.
\]
\paragraph{Search distance.}
\[
  D = \frac{1}{2}\sum_{r=1}^{2} \frac{1}{|\mathcal{A}_r|\,B_r}
      \sum_{t\in\mathcal{A}_r}\sum_{b=1}^{B_r}
      \bigl|\,T_r(t,b) - \bigl(R_r(t,b) + e_r(b) + o_r\bigr)\bigr|.
\]
In words: the average number of decibels by which the render misses the target
in each band of each active frame, \emph{after} the render has been given the
best smooth EQ (within $\pm18\dB$) that matches the long-term spectra.

\paragraph{Benchmark score} (\code{plain_distance}): the same formula without
$e_r$ and $o_r$, and always with the 80\dB{} floor. It measures how far apart
the two signals sound, colour included.

\begin{keyidea}[: why forgive an EQ during the search]
The stand-in dry signal has the right notes but the wrong guitar, and the target
went through a cabinet, a microphone and a mix EQ that the template chain does
not model. Both are mostly \emph{linear} differences, which a filter fixes
afterwards (Section~\ref{sec:dsp}). Forgiving them lets the search spend its
effort on what only it can fix: input level, compression, drive, amp gain and
time-based effects. Distortion is \emph{not} forgiven: clipping changes the
spectrum frame by frame, not by a fixed EQ (a unit test checks this:
\code{distortion_is_not_forgiven}).
\end{keyidea}

\subsection{The template chain}

\file{tone/template.rs}. The search moves a vector $x\in[0,1]^n$; each
coordinate maps to a pedal parameter, linearly or on a log scale
(Table~\ref{tab:template}). Order: compressor $\to$ distortion $\to$ amp $\to$
delay $\to$ reverb. A pedal whose amount (compressor amount, drive mix, delay
mix, reverb mix) is below 0.04 is left out of the chain, so the search can
switch pedals off.

\begin{table}[h]
\centering
\caption{Searched parameters. The first ten are in the committed version
\done; the two delay parameters are in the working tree \wip.}
\label{tab:template}
\small
\begin{tabular}{@{}rllll@{}}
\toprule
\# & Parameter & Range & Scale & Maps to \\
\midrule
0 & input level (RMS) & 0.01--0.4 & log & level the dry signal is played at \\
1 & compressor amount & 0--1 & lin & ratio $1 + 7\cdot$amount (up to 8:1) \\
2 & compressor threshold & 0.02--0.5 & log & threshold \\
3 & drive mix & 0--1 & lin & Distortion mix \\
4 & drive & 0.5--10 & log & Distortion drive \\
5 & drive tone & 0.15--1 & lin & Distortion tone \\
6 & amp gain & 0.5--5 & log & Amp gain (voicing searched separately) \\
7 & reverb mix & 0--0.5 & lin & Reverb mix \\
8 & reverb room & 0--0.95 & lin & Reverb room size \\
9 & reverb damping & 0--0.95 & lin & Reverb damping \\
10 & delay mix & 0--0.5 & lin & Delay mix (time measured, not searched) \\
11 & delay feedback & 0--0.7 & lin & Delay feedback \\
\bottomrule
\end{tabular}
\end{table}

Fixed values: compressor attack 10\ms, release 150\ms, make-up gain 1;
distortion level 1, bias 0; amp tone stack flat (bass, mid, treble 0), because
it acts after the clipping and is therefore linear and left to the match IR;
delay with one tap and damping 0.3. The neutral start is input RMS 0.1,
compressor amount 0.2, drive mix 0.3, reverb mix 0.15, delay mix and feedback 0.3,
everything else 0.5.

\subsection{The search}

\file{tone/cmaes.rs} implements separable CMA-ES~\cite{ros2008sepcmaes}
(Section~\ref{sec:search}) with a deterministic random generator (SplitMix64), so
the same seed gives the same result. Candidates outside the unit cube are
evaluated at the nearest point inside, plus a penalty of 10 times the squared
distance to the cube. Each generation's candidates are rendered in parallel.
Settings: population 16, 30 generations, initial step size 0.3, seed
incremented per voicing. The amp voicing is discrete, so there is one search per
voicing.

\subsection{Measuring the delay time}
\label{sec:impl-delay}

\file{tone/analysis.rs} \wip.
\begin{keyidea}[: echoes repeat onsets]
A delay repeats every onset of the target after the delay time. The stand-in dry
signal has the same notes without echoes. So a time lag at which the target's
onset pattern correlates with itself much more than the dry signal's does is the
delay time; rhythmic repetition in the playing shows up in both and cancels.
\end{keyidea}
Both signals are turned into an onset envelope (positive spectral flux of the log
magnitude, 1024-point FFT, one value every 5\ms). For each lag from 60\ms{} to
1.2\,s the code computes the normalised autocorrelation of the target envelope
minus that of the dry envelope (the \emph{excess}). If the largest excess is
below 0.12 there is no echo and no delay pedal. Otherwise it takes the
\emph{shortest} lag whose excess reaches 70\% of the best (a feedback delay also
correlates at two and three times its time), refined to the local maximum. The
benchmark counts a delay time within 5\% of the truth as right.

\subsection{The match filter}
\label{sec:impl-match}

\file{tone/matching.rs} \done. The EQ $e(b)$ of the finer-frequency resolution
(40 bands) is turned into an impulse response:
\begin{enumerate}
\item Interpolate the EQ on a log-frequency axis at every bin of an 8192-point
  FFT (flat beyond the outer bands).
\item Take the natural log of the magnitude and its inverse FFT: the real
  cepstrum.
\item Fold it: keep coefficient 0 and $N/2$, double the causal half, zero the
  rest. This is the cepstrum of the minimum-phase filter with the same
  magnitude~\cite{oppenheim2010dtsp}.
\item FFT, complex exponential, inverse FFT: the minimum-phase IR. Keep the
  first 2048 samples (43\ms) and fade the last quarter with a half-cosine so the
  cut does not ring.
\end{enumerate}
The IR is written as a WAV file and loaded by a Cabinet pedal (mix 1) appended at
the end of the chain (\code{Candidate::chain_with_ir}).

\subsection{Saving the result and the UI}
\label{sec:impl-save}

\wip{} (working tree, not committed at the time of writing).
\code{save_preset} writes the IR to \file{<data dir>/guitar_virtual_amp/tone_match/}
and a user preset named ``\emph{title} -- \emph{track} (timbre estimado)'' with
\code{song} and \code{artist} set and the tags \code{timbre-estimado} and
\code{faixa:<track>}; estimating again overwrites it.
\code{shared::tab::tone::preset_fits_track} keeps a preset tagged for one track
from being suggested for another. In the Transcribe form, ``Estimar o timbre
(experimental)'' (\code{AppSettings::transcription_tone}, off by default, needs
separation) adds a last stage to \code{transcribe::run} that estimates each
guitar part from its own audio (its pan region of the stem) and notes; the form
says it takes about 30\,s per guitar.

% !TeX root = ../main.tex
\section{How well it works}
\label{sec:results}

\subsection{The benchmark}

\code{eval_tone} (an ignored test in \file{tone/eval.rs}) uses EGDB test items:
each has the DI of a performance, renders of the same DI through three
amplifier simulations (Marshall, Mesa, Plexi: equipment \emph{outside} the app's
chain) and labelled notes. Per item there are four kinds of target:
\begin{description}[style=nextline,leftmargin=1.2cm]
\item[chain] The DI through a random template chain followed by a random
  cabinet-like EQ (a high-pass near 80\Hz, a roll-off above 3--7\kHz, six
  smooth bumps of up to $\pm4\dB$). The true settings are known, so parameter
  error can be measured. Random settings: a random voicing, each parameter
  uniform in 0.1--0.9, and the compressor, drive and reverb each switched off
  with probability 0.5, 0.4 and 0.5. In the working tree half of these targets
  also get a delay of 150--700\ms{} \wip.
\item[Marshall, Mesa, Plexi] The EGDB renders: can the app's chain imitate gear
  it does not contain?
\end{description}

Each target is estimated by five methods, all ending in the same kind of chain:
\begin{enumerate}
\item \textbf{dry DI}: no processing at all (a floor to beat).
\item \textbf{best factory preset}: the closest of the 20 presets, as is.
\item \textbf{preset + match}: that preset plus the match IR.
\item \textbf{search + match, from notes}: the full estimator, with the stand-in
  dry signal from the plucked-string model playing the \emph{labelled} notes.
  This is the real use case, except that the notes are perfect.
\item \textbf{search + match, true DI}: the full estimator given the real DI.
  An upper bound: what the search achieves when the dry signal is right.
\end{enumerate}

\begin{keyidea}[: how each estimate is scored]
Whatever dry signal the estimator used, the \emph{true} DI is rendered through
its estimated chain and compared with the target by the plain log-band distance
(no EQ allowance, Section~\ref{sec:app-distance}). This asks the question the
user cares about: if I play this performance through the proposed preset, how
far is it from the record? It also makes the five methods comparable.
\end{keyidea}

\subsection{Results}

Committed run (\code{7db5015}): 12 EGDB items, 12\,s excerpts, about 15\,s per
estimate on 28 threads \measured.

\begin{table}[h]
\centering
\caption{Mean log-band distance (dB, lower is better) between the true DI
rendered through each estimate and the target, 12 EGDB items \measured.}
\label{tab:tone-results}
\footnotesize
\setlength{\tabcolsep}{4pt}
\begin{tabular}{@{}lccccc@{}}
\toprule
Target & Dry DI & Best preset & Preset + match & Search, notes & Search, true DI \\
\midrule
Random app chain + cab EQ & 6.60 & 9.29 & 6.64 & \textbf{5.30} & 0.66 \\
EGDB Marshall & 7.52 & 7.55 & 6.84 & \textbf{6.49} & 3.11 \\
EGDB Mesa & 7.94 & 8.90 & 6.52 & \textbf{5.83} & 2.67 \\
EGDB Plexi & 7.61 & 8.16 & 7.63 & \textbf{6.26} & 2.32 \\
\bottomrule
\end{tabular}
\end{table}

\paragraph{What the table says.}
\begin{itemize}
\item \textbf{The estimator works, modestly.} Search from notes beats every
  baseline on average in all four groups, and in 7--8 of the 12 items per group.
  The gain over the best baseline is 0.35--1.35\dB.
\item \textbf{The match filter matters.} A factory preset alone is often
  \emph{worse} than no processing (9.29 against 6.60 on chain targets), because
  its colour is wrong; with the match IR it gets close to or better than the dry
  DI.
\item \textbf{The search itself is good; the stand-in dry signal is the
  bottleneck.} With the true DI the search gets within 0.66\dB{} of targets made
  with the app's own chain, essentially solving them, and to 2.3--3.1\dB{} of
  real amp simulations. From the notes it stays at 5.3--6.5\dB. The difference
  (3.2--4.6\dB) is the cost of the plucked-string stand-in, which is still
  9.9\dB{} from real DI even with an EQ allowed.
\item \textbf{Settings are not recovered, and that is expected.} On chain targets
  the mean parameter error is 0.41 on the 0--1 scale and the amp voicing is right
  in 7 of 12 items. Several settings sound alike (Section~\ref{sec:comparison}),
  so the search finds \emph{a} chain that sounds close, not \emph{the} chain.
\end{itemize}

\begin{worked}[: what a decibel of distance means]
The score is an average absolute error per band and frame. 6\dB{} means each
band of each frame is, on average, off by about a factor of two in amplitude
(20\,log$_{10}$\,2 $\approx$ 6\dB). It is a coarse measure: it cannot tell a
small error everywhere from a large error in a few bands, and it has not been
checked against listeners for these sounds.
\end{worked}

\begin{caution}[: what this benchmark does not yet test]
\begin{itemize}
\item \textbf{Perfect notes}: the stand-in plays labelled notes. With
  transcribed notes, missed and extra notes will add error \expected.
\item \textbf{Isolated guitar}: no target is a separated stem. The working tree
  adds separated-stem targets (\code{GVA_TONE_STEMS}), scored against the clean
  render they came from \wip.
\item \textbf{Small sample}: 12 items. Differences of a few tenths of a dB
  between methods are within item-to-item variation.
\item \textbf{The score is not validated}: no listening test yet. The benchmark
  writes WAVs (\code{GVA_TONE_OUT}) for listening.
\item \textbf{Only simulated amps} for the gear-outside-the-chain case; the
  working tree adds Guitar-TECHS miked real amplifiers \wip.
\end{itemize}
\end{caution}

\subsection{Being measured in the working tree}

At the time of writing, the benchmark in the working tree \wip{} adds: the
sampler stand-in as a sixth method (\code{GVA_TONE_SAMPLES}); delay-time accuracy
(right within 5\% when the target has a delay, ``no delay'' when it has none);
separated-stem and Guitar-TECHS miked-amp targets; a higher search floor
(\code{GVA_TONE_FLOOR}) so the search ignores quiet bands where a stand-in
differs most from a real DI; and \code{GVA_TONE_PLUCK=0} to skip the pluck
stand-in. No numbers from it are recorded yet; when they are, this section and
Table~\ref{tab:tone-results} should be updated.

% !TeX root = ../main.tex
\section{Next steps}
\label{sec:next}

\subsection{Order of work}

The results point at the dry signal first: the search is close to optimal when
it has the right one.

\begin{nextsteps}
\begin{enumerate}
\item \textbf{A better dry signal} \wip: the sampler of real DI notes; later a
  learned DI renderer (in the manner of GuitarFlow~\cite{loth2025guitarflow},
  trained on redistributable DI), an effect remover trained on the app's own
  renders~\cite{lee2024distortion,yeh2026egvae}, and the player's own DI through a
  pre-chain tap in practice mode.
\item \textbf{Separated stems} in the benchmark \wip, then the transcription
  pipeline's own notes instead of labels.
\item \textbf{Time-based effects from notes and tempo}: delay \wip; reverb decay
  from rests; modulation rate from sustained notes \planned.
\item \textbf{A learned initial guess} trained on separated renders of the
  template chain, returning several candidates \planned.
\item \textbf{Listening tests} (MUSHRA via webMUSHRA) to check that the distance
  agrees with ears \planned.
\item \textbf{The UI}: save as a preset tagged with the song and track \wip;
  compare the preset with the stem played in sync by the tab player \planned.
\item \textbf{Captures}: the user's NAM files as amp candidates \planned.
\end{enumerate}
\end{nextsteps}

\subsection{The full evaluation protocol}

The benchmark is meant to grow into four sets:
\begin{enumerate}
\item \textbf{In-chain synthetic}: held-out DI (Guitar-TECHS P3, GuitarSet
  players 03--05) through random template chains, isolated and over backing tracks
  after separation; parameter error and audio distances. The difference between
  isolated and separated inputs measures the cost of separation.
\item \textbf{Equipment outside the chain}: EGDB and EGDB-PG renders, Guitar-TECHS
  real amplifiers, and for research only GOAT and ToneTwist devices; the true DI
  is rendered through the estimate and compared with the real output.
\item \textbf{Real productions}: MoisesDB songs with isolated guitar tracks;
  estimates from the separated stem against estimates from the isolated track,
  each scored against the isolated track.
\item \textbf{Listening}: MUSHRA on sets 1 and 2, where the same performance
  exists on both sides; tone-similarity scores on set 3; correlation of the
  objective metrics with the ratings.
\end{enumerate}
Baselines throughout: the current metadata-based suggestion, the nearest factory
preset, that preset plus the match, the full pipeline, and the search with the
true DI as upper bound.

\subsection{Conclusion}

Each part of guitar tone estimation has working methods under easier conditions:
effects are recognised on synthetic isolated guitar, amplifiers are captured from
paired audio and, less accurately, from unpaired audio, tone-transfer models
reproduce whole chains as audio, and inference-time search drives arbitrary,
non-differentiable chains. What was missing is the combination the commercial
case needs: an editable chain, estimated from a separated stem without the dry
signal. The app now has a first version of that combination, using its own
renderer as the forward model, a gradient-free search, a final linear match,
and one idea the literature had not used: a dry signal synthesised from the
transcribed notes. The first benchmark shows that the search is not the weak
part; the stand-in dry signal is. Improving it, and measuring on separated stems
and with listeners, are the next steps.


\appendix
\part*{Appendices}
\addcontentsline{toc}{part}{Appendices}
% !TeX root = ../main.tex
\section{Glossary}
\label{app:glossary}

\begin{description}[style=unboxed,leftmargin=0cm,itemsep=3pt]
\item[Analysis by synthesis] Estimating settings by rendering candidates through
  the forward model and comparing them with the target.
\item[Cabinet IR] The impulse response of a speaker cabinet and microphone,
  applied by convolution.
\item[Capture] A neural model of one device at one setting, learned from
  dry/wet pairs (NAM, Kemper, TONEX).
\item[Cepstrum] Inverse Fourier transform of the log magnitude spectrum; used to
  build minimum-phase filters.
\item[CMA-ES] Covariance Matrix Adaptation Evolution Strategy, a gradient-free
  optimiser (\S\ref{sec:search}); \emph{separable}: diagonal covariance only.
\item[Convolution] Applying an impulse response to a signal.
\item[DDSP] Differentiable digital signal processing: DSP blocks written so that
  gradients flow through them.
\item[DI] Direct input: the dry guitar signal before any processing.
\item[Dry / wet] Unprocessed / processed signal.
\item[EQ allowance] The smooth, capped EQ the app's search distance forgives
  (\S\ref{sec:app-distance}).
\item[ESR] Error-to-signal ratio, the usual training loss of amp models.
\item[FAD] Fréchet Audio Distance between the embedding distributions of two
  sets of clips.
\item[Forward / inverse problem] Settings $\to$ sound / sound $\to$ settings.
\item[Grey-box model] A model with the circuit's structure and learned values.
\item[Impulse response (IR)] The output of an LTI system for a single click;
  describes the system completely.
\item[LFO] Low-frequency oscillator driving a modulation effect.
\item[Long-term average spectrum] A signal's spectrum averaged over a long
  stretch; differences between two of them give a matching EQ.
\item[LTI] Linear and time-invariant.
\item[Minimum phase] Of all filters with a given magnitude response, the one
  with the least delay.
\item[MR-STFT] Multi-resolution STFT distance.
\item[MUSHRA] A standard listening-test method with hidden reference and anchor.
\item[NAM] Neural Amp Modeler, the open capture format the app's NAM pedal plays.
\item[Pseudo-DI / stand-in dry signal] A dry signal not recorded with the target
  but made for it (here: from the notes).
\item[RMS] Root mean square, a measure of level.
\item[SI-SDR] Scale-invariant signal-to-distortion ratio, in dB.
\item[Stem] One instrument's audio separated from a mix.
\item[Template chain] The fixed pedal chain the app's search fits
  (Table~\ref{tab:template}).
\item[Tone stack] The bass/middle/treble filter of an amplifier.
\item[Unpaired / wet only] Dry audio from another performance is available /
  no dry audio at all.
\end{description}

% !TeX root = ../main.tex
\section{Code map and how to reproduce the numbers}
\label{app:code}

\subsection{Source files}

\begin{center}
\small
\begin{tabularx}{\textwidth}{@{}>{\raggedright}p{5.4cm}X@{}}
\toprule
File (under \file{guitar_core/src/}) & Role and section \\
\midrule
\file{tone/mod.rs} & \code{estimate}, baseline, per-voicing search, \code{save_preset} (\S\ref{sec:impl}) \\
\file{tone/pluck.rs} & plucked-string stand-in and \code{DI_EQ} (\S\ref{sec:impl-dry}) \\
\file{tone/sampler.rs} & sampler stand-in from real DI notes (\S\ref{sec:impl-dry}, \wip) \\
\file{tone/features.rs} & log-band distance with EQ allowance; \code{plain_distance} (\S\ref{sec:app-distance}) \\
\file{tone/template.rs} & template chain and parameter ranges (Table~\ref{tab:template}) \\
\file{tone/cmaes.rs} & separable CMA-ES (\S\ref{sec:search}) \\
\file{tone/analysis.rs} & delay time from echoes of onsets (\S\ref{sec:impl-delay}, \wip) \\
\file{tone/matching.rs} & minimum-phase match IR, FFT convolution (\S\ref{sec:impl-match}) \\
\file{tone/eval.rs} & \code{eval_tone}, \code{fit_pluck}, \code{pluck_vs_di} (\S\ref{sec:results}) \\
\file{backend/offline.rs} & \code{render_at}: the chain offline at a given input level \\
\file{shared/tab/tone.rs} & tone suggestions; estimated-preset tags (\wip) \\
\bottomrule
\end{tabularx}
\end{center}

\subsection{Commands}

\begin{verbatim}
GVA_TONE_MANIFEST=~/.cache/gva_test/prepared/eval_egdb.tsv \
  cargo test --release --workspace --exclude guitar_android eval_tone \
  -- --ignored --nocapture
\end{verbatim}

\begin{center}
\small
\begin{tabularx}{\textwidth}{@{}>{\raggedright}p{4.2cm}X@{}}
\toprule
Variable & Effect \\
\midrule
\code{GVA_TONE_LIMIT} & number of items (default 8; the recorded run used 12) \\
\code{GVA_TONE_GENS} & CMA-ES generations (default 30) \\
\code{GVA_TONE_SECS} & excerpt length in seconds (default 12) \\
\code{GVA_TONE_ORACLE=0} & skip the true-DI search \\
\code{GVA_TONE_OUT=<dir>} & write WAVs (target, estimate, preset) for listening \\
\code{GVA_TONE_STEMS=<tsv>} & separated-stem targets (\wip) \\
\code{GVA_TONE_SAMPLES=<di:notes,...>} & sampler stand-in (\wip) \\
\code{GVA_TONE_FLOOR} & search floor in dB (\wip) \\
\code{GVA_TONE_PLUCK=0} & skip the pluck stand-in (\wip) \\
\bottomrule
\end{tabularx}
\end{center}

\code{pluck_vs_di} and \code{fit_pluck} take
\cmd{GVA_PLUCK_ITEMS=<di.wav>:<notes.txt>,...} (Guitar-TECHS DI only, never EGDB).


\bibliographystyle{plainnat}
\bibliography{references}

\end{document}
